% Pandoc template -> Elsevier elsarticle, for CAOR-D-26-01812.
% Driven by scripts/build_submission.py.
%
% CAOR Guide for Authors (ScienceDirect, accessed 2026-08-19):
%   - single-anonymized review
%   - Harvard (author-year) bibliography
%   - highlights: 3-5 bullets, <=85 characters each (also in paper/highlights.txt)
%   - abstract <= 250 words
%   - keywords: 1-7
%   - data statement required (Option C)
%   - CRediT required
%   - no generative-AI declaration in this source (authors will add at submission)
%
% The `review` option is Elsevier's double-spaced single-column review layout.
% Swap to `final,5p,times,twocolumn` for a preprint that looks like the article.
\documentclass[review,3p,times]{elsarticle}

\usepackage[T1]{fontenc}
\usepackage[utf8]{inputenc}
\usepackage{amsmath,amssymb,amsthm}
\usepackage{booktabs}
\usepackage{longtable}
\usepackage{array}
\usepackage{calc}
\usepackage{graphicx}
\usepackage{xcolor}
\usepackage{hyperref}
\hypersetup{colorlinks=true,linkcolor=black,citecolor=black,urlcolor=blue}

\usepackage{adjustbox}
\providecommand{\tightlist}{%
  \setlength{\itemsep}{0pt}\setlength{\parskip}{0pt}}
\providecommand{\pandocbounded}[1]{#1}
\renewcommand{\pandocbounded}[1]{\adjustbox{max width=\linewidth,center}{#1}}

\DeclareUnicodeCharacter{0107}{\'{c}}
\DeclareUnicodeCharacter{0161}{\v{s}}
\DeclareUnicodeCharacter{2265}{\ensuremath{\geq}}
\DeclareUnicodeCharacter{2713}{\ensuremath{\checkmark}}
\DeclareUnicodeCharacter{00D7}{\ensuremath{\times}}
\DeclareUnicodeCharacter{2013}{--}
\DeclareUnicodeCharacter{2014}{---}
\DeclareUnicodeCharacter{2212}{-}
\newcounter{none}

% pandoc --citeproc already inlines Harvard (author-year) citations.
% Do not wrap the list in thebibliography: natbib then demands its labels.
\newcommand{\citeproctext}{}
\newenvironment{CSLReferences}[2]{%
  \par\addvspace{\baselineskip}%
  \renewcommand{\bibitem}[2][]{\par\addvspace{0.35\baselineskip}\noindent\hangindent=1.5em\hangafter=1}%
}{\par}

\newtheorem{proposition}{Sketch}
\newtheorem{assumption}{Assumption}
\theoremstyle{definition}
\newtheorem{definition}{Definition}
\newtheorem{remark}{Remark}


\journal{Computers \& Operations Research}

\begin{document}

\begin{frontmatter}

\title{Offline truncated-rollout labels did not recover the online
teacher on one warehouse simulator}

\author[rvce]{Vittal Mukunda\corref{cor1}}
\ead{vittalmukunda.im24@rvce.edu.in}
\author[rvce]{Atharva Somani}
\ead{atharvasomani.im24@rvce.edu.in}
\author[rvce]{Pranjal Malaiya}
\ead{pranjalmalaiya.im24@rvce.edu.in}

\affiliation[rvce]{organization={Department of Industrial Engineering and Management, R. V. College of Engineering},
                   addressline={Mysore Road},
                   city={Bengaluru},
                   postcode={560059},
                   country={India}}

\cortext[cor1]{Corresponding author.}

\begin{abstract}
Warehouse dispatching under due dates and product expiry is usually left
to a priority rule. A selector can pick the rule from a queue summary by
rolling every candidate forward from logged states and fitting a
classifier to the cost vectors. That construction is not new. This paper
asks whether the distilled classifier recovers the online truncated
lookahead that produced its labels, on one warehouse simulator.

Environment, seeds and objective are held fixed on one warehouse
simulator with two priced clocks. The distilled ranker (DAHS) is a
gradient-boosted tree. Comparators are the online teachers, static
rules, a one-step snapshot, fitted Q-iteration, and a neural PPO policy
that is not a matched simulation budget. Live evaluation admits leftover
arrivals at shift end (mean arrived orders 791). Production labels
omitted that admit.

On 50 held-out shifts DAHS has mean composite cost \(J=382\). Online
lookahead is cheaper (\(357\)--\(363\)). Always-COVERT costs \(455\).
Fitted Q costs \(398\); that 4\% gap is not a result about value
learning. A one-step snapshot matches DAHS on the paired mean interval.
A two-step label is cheaper than the deployed four-step ranker on the
mean (\(-9.11\), \([-17.61,-1.35]\); 21/20/9 wins-losses-ties, median
0). On a 12-cell grid the one-step teacher wins eight cells, Always-EEDD
wins four, and DAHS wins none. Distillation takes \(4.24\) ms against
\(670\) ms of lookahead, which is \(0.07\%\) of a 15-minute review. If a
resettable simulator exists, run the teacher. If it does not, these
labels cannot be generated.
\end{abstract}

\begin{highlights}
\item On one simulator, distillation misses the teacher (J 382 vs 357).
\item The distilled selector wins 0 of 12 robustness cells; the teacher wins 8.
\item A two-step label is cheaper on the mean than the four-step ranker.
\item 670 ms lookahead is 0.07\% of a 15-minute review interval.
\item Soft labels, extra samples, and 33 of 38 features are null on cost.
\end{highlights}

\begin{keyword}
dynamic dispatching \sep selection
hyper-heuristics \sep rollout \sep approximate policy
iteration \sep warehouse operations \sep perishable inventory
\end{keyword}

\end{frontmatter}

\section{Introduction}\label{introduction}

Order dispatching on a warehouse floor is a sequential decision problem
under uncertainty: orders arrive stochastically, each carries a due date
and possibly a product expiry, and a small pool of pickers must be
assigned work so as to minimise late and spoiled shipments. In practice
the decision is delegated to a \emph{dispatching rule}
(first-in-first-out, earliest-deadline-first, weighted
shortest-processing-time, and the like) because rules are transparent,
fast, and require no training. No single rule dominates: the rule that
minimises lateness under light load is not the rule that does so when
the queue is saturated or when a burst of perishable orders arrives. A
controller that \emph{selects} the rule appropriate to the current
state, a \emph{selection hyper-heuristic} (Drake et al. 2020; Dokeroglu
et al. 2024), can in principle capture the envelope of the pool without
abandoning the operational advantages of rules.

The open question is how to \emph{learn} the selector. Deep
reinforcement learning (Mahmoudinazlou et al. 2025; Chen et al. 2024) is
sample-hungry, and it is unstable on problems where the per-state
advantage of one action over another is small relative to the return
variance, which is the usual case in rule selection. Imitation of an
expert dispatcher (Han and Jung 2025) needs an expert. The training
signal we study has neither problem: it measures the counterfactual cost
of every rule directly.

The signal is a rollout (Bertsekas 2020). Fix a state, run each
candidate rule forward for a short horizon, and record the cost it
incurs. Rollouts can be run online at each decision, or offline once
over a corpus of states. For each state in a corpus of simulated shifts
we roll out every rule, obtain a per-rule cost vector, and fit a
supervised ranker to it. At deployment the ranker is a single forward
pass. We retain the cost margin between rules through a soft,
tempered-softmax label by default; an ablation (Section 6.8) shows the
soft form is not essential.

Simulating a rule pool offline and fitting a classifier to the result is
not a new training mechanism. It is related to rollout classification
policy iteration and to multi-pass rule selection; Section 2 places the
method inside both. The construction here is a single supervised fit to
truncated-rollout labels, not an iterated policy-iteration loop, so we
have not implemented RCPI.

The question this paper answers is narrower, and empirical. On one
warehouse shift corpus, with environment, seeds and objective held
fixed, does the distilled ranker recover the online truncated lookahead
that generated its labels? Secondary comparators vary the training
signal: a bootstrapped state-action value fitted from the same logged
transitions (same gradient-boosted-tree class as the ranker), and a
neural PPO policy that is not the same approximator class and not a
matched simulation-step budget. Section 6.10 reports that comparison;
Section 6.9 reports the PPO hyperparameter grid. The warehouse model
carries two priced clocks, customer due date and product expiry, and
Section 3.5 measures whether the second binds at a 15-minute review. The
name DAHS expands to Disruption-Aware Heuristic Scheduling. The
simulated warehouse is stationary: Poisson arrivals, i.i.d. processing
times, no breakdowns, no cancellations. Section 6.11 perturbs parameters
between labelling and evaluation; that is model misspecification, not a
disruption process.

The empirical answer is no. On the default operating point DAHS beats
every static dispatching rule we screened, with Always-COVERT as the
static to beat. It does not beat the online teachers. On a 12-cell
robustness grid five methods are frozen: the one-step teacher wins eight
cells, Always-EEDD wins four, and DAHS wins none. \(M\) from 1 to 40
sits in a \(0.7\%\) band on deployed \(J\). Under the confirmatory
paired interval, \(\tau=1\) is a null against deployed \(\tau=4\);
\(\tau=2\) is cheaper than \(\tau=4\) on the mean, with 21 wins, 20
losses, 9 ties and median difference 0. The deployed model remains
\(\tau=4\) because that is the frozen labelling budget. Distillation is
\(158\) times faster than the rolling-horizon teacher (\(4.24\) ms
against \(670\) ms). At a 15-minute review, \(670\) ms is \(0.07\%\) of
the epoch. If a resettable simulator is available and \(670\) ms is
affordable, the teacher is the cheaper policy. If it is not,
truncated-rollout labels of this kind cannot be generated on this
corpus.

The rest of the paper is the protocol, the scoreboard, and the
limitations of a single simulator.

\subsection{Terminology and notation}\label{terminology-and-notation}

Terms such as ``corpus of simulated shifts'', ``held-out shifts'',
``SLA-breach rate'' and ``snapshot-trained ranker'' are defined here, on
first use in the text, or both.

{\def\LTcaptype{none} % do not increment counter
\begin{longtable}[]{@{}
  >{\raggedright\arraybackslash}p{(\linewidth - 2\tabcolsep) * \real{0.5000}}
  >{\raggedright\arraybackslash}p{(\linewidth - 2\tabcolsep) * \real{0.5000}}@{}}
\toprule\noalign{}
\begin{minipage}[b]{\linewidth}\raggedright
Term
\end{minipage} & \begin{minipage}[b]{\linewidth}\raggedright
Meaning
\end{minipage} \\
\midrule\noalign{}
\endhead
\bottomrule\noalign{}
\endlastfoot
\textbf{decision epoch} & The boundary of a review interval, where the
controller acts. There are \(N\) per shift. \\
\textbf{dispatching rule} & A function that orders the waiting queue.
Pickers are then assigned down that order. FIFO, EDD and the rest are
dispatching rules. \\
\textbf{selection hyper-heuristic} & A controller that chooses
\emph{which dispatching rule to apply}, as a function of the current
state, rather than choosing an assignment directly. \\
\textbf{corpus of simulated shifts} & A set of shifts generated from
distinct random seeds, used as data. Split into three disjoint blocks:
\textbf{training} (fits the selector), \textbf{calibration} (fits rule
parameters such as ATC's look-ahead scale), and \textbf{test}. \\
\textbf{held-out shifts} & The test block used for confirmatory tables.
Ranker grid, isotonic split, ATC/COVERT \(k\) and regime \(K\) are fit
on training or calibration, not on test. SHAP values and the
\(T_{\min}\) dwell sweep are computed on this block and are not used to
select the deployed model. A \texttt{top5\_features} retrain that read
test-set SHAP is reported as a diagnostic in Section 6.8, not in the
ablation table. The PPO sensitivity grid in Section 6.9 also read this
block; the main comparison keeps the pre-declared \texttt{ppo\_fair}
row. \\
\textbf{SLA} & Service-level agreement --- the contractual delivery
commitment. An order's SLA due time \(d_o\) is when it must ship. \\
\textbf{SLA-breach rate} & The fraction of orders shipped after \(d_o\).
Two denominators are possible and the choice matters (Section 3.3): over
\emph{arrived} orders, or over \emph{completed} orders only. We always
state which. \\
\textbf{service-failure rate} & Our primary reported KPI component: the
fraction of \emph{arrived} orders that either ship late or spoil,
whether or not they were ever dispatched. \\
\textbf{rollout} & Simulating forward from a state under a fixed rule,
to measure what that rule costs. \\
\textbf{truncated rollout} & A rollout stopped after \(\tau\) intervals
instead of running to the end of the shift. \\
\textbf{continuation} & One independently sampled future used for a
rollout. The label averages over \(M\) of them. \\
\textbf{ranker} & The fitted classifier that maps an observation to a
distribution over rules. A gradient-boosted decision-tree ensemble
here. \\
\textbf{snapshot-trained ranker} & The ranker fitted to labels from a
\emph{one-interval} rollout, i.e.~\(\tau = 1\). Used as an ablation to
isolate the value of looking further ahead. \\
\textbf{soft label} & The training target expressed as a distribution
over rules rather than a single winner, obtained from the cost vector by
a tempered softmax. \\
\textbf{regime} & A cluster of operating conditions, fitted by a
Gaussian mixture over training observations. Membership probabilities
are appended to the observation. \\
\textbf{switching controller} & The deployment wrapper that enforces a
minimum \textbf{dwell} (hold a rule for \(T_{\min}\) epochs) and an
\textbf{entropy gate} (permit an early switch when the ranker is
confident). \\
\textbf{ablation} & Removing or altering one component and re-measuring,
to establish what that component contributes. \\
\textbf{DAHS} & Disruption-Aware Heuristic Scheduling, the selection
hyper-heuristic studied here. The name is the code name. The environment
in Section 3 does not model disruptions. \\
\end{longtable}
}

\textbf{Notation.}

{\def\LTcaptype{none} % do not increment counter
\begin{longtable}[]{@{}
  >{\raggedright\arraybackslash}p{(\linewidth - 2\tabcolsep) * \real{0.5000}}
  >{\raggedright\arraybackslash}p{(\linewidth - 2\tabcolsep) * \real{0.5000}}@{}}
\toprule\noalign{}
\begin{minipage}[b]{\linewidth}\raggedright
Symbol
\end{minipage} & \begin{minipage}[b]{\linewidth}\raggedright
Meaning
\end{minipage} \\
\midrule\noalign{}
\endhead
\bottomrule\noalign{}
\endlastfoot
\(t\), \(N\), \(L\), \(T\) & epoch index; epochs per shift (32);
interval length (15 min); shift length (480 min) \\
\(S_t\) & the true state: queue with full per-order attributes, picker
availability, clock \\
\(\phi(\cdot)\), \(x_t\) & the feature map, and the observation
\(x_t = \phi(S_t)\) the controller actually sees \\
\(u_t\) & the decision at epoch \(t\) --- a rule from the pool \\
\(W_{t+1}\) & exogenous information: orders arriving in \((t, t+1]\) and
their attributes \\
\(S^M(\cdot)\) & the transition function (Section 3.4) \\
\(\mathcal{H}\), \(h\) & the rule pool and a member of it \\
\(H_t\) & intervals remaining in the shift at epoch \(t\) (distinct from
\(\mathcal{H}\)) \\
\(a_o, p_o, d_o, x_o, w_o\) & order \(o\)'s arrival, processing time,
customer deadline, product expiry, economic weight \\
\(f_o\) & completion time of order \(o\) \\
\(J\); \(W_b, W_t, W_s, W_h\) & the composite objective and its weights:
breach, tardiness, spoilage, holding \\
\(\tau\), \(M\), \(\beta\) & rollout horizon; continuations per rollout;
dimensionless softmax-temperature multiplier
(\(T(s)=\beta\hat\sigma(s)\)) \\
\(T(s)\) & per-state softmax temperature, in cost units \\
\(\varepsilon\) & per-step model error in total variation, used only as
a description of the misspecification grid in Section 6.11 \\
\end{longtable}
}

The remainder of the paper is organised as follows. Section 2 reviews
related work. Section 3 defines the dispatching problem and the
simulator. Section 4 presents DAHS. Section 5 describes the experimental
protocol. Section 6 reports results. Section 7 is the practitioner
reading; Section 8 lists limitations; Section 9 concludes.

\section{Related Work}\label{related-work}

\subsection{Scope: which order-picking decision this paper
addresses}\label{scope-which-order-picking-decision-this-paper-addresses}

Order picking decomposes into a set of interacting planning problems ---
storage assignment, zoning, order batching, picker routing, and order
release and dispatching (Koster et al. 2007; Gils et al. 2018). This
paper addresses \textbf{order release and dispatching only}: at a fixed
review epoch, which of the waiting orders are released to which picker,
and in what order. Three neighbouring problems are treated as exogenous
and are held fixed throughout, and because that restriction shapes every
result we report, we state it explicitly rather than leaving it to be
inferred.

\textbf{Storage assignment and layout} are fixed. They are themselves
substantial design problems (Roodbergen and Vis 2006), and they enter
this model only through the distribution of order processing times.

\textbf{Batching} is fixed at one order per pick tour. The controller
never decides which orders to combine; it decides only which order a
free picker starts next.

\textbf{Routing and travel-time estimation} are exogenous. The
processing time \(p_o\) is a single realised service duration that
aggregates travel, search and pick time, in the three-point form
standard when only time-standard data are available (Tompkins et al.
2010).

This restriction is deliberate for three reasons. First, it matches the
planning horizons of the decisions involved: layout, zoning and batching
policy are typically fixed over horizons far longer than a shift, so
within-shift adaptive control acts on dispatching alone (Gils et al.
2018). Second, rule selection is only a well-posed question once the
downstream problems are fixed; if the action space also ranged over
batch composition or route construction, the comparison between
dispatching rules would confound the rule with the batching and routing
policy it happens to be paired with. Third, travel-time estimation is a
substantial modelling problem in its own right, and embedding an
estimator inside the environment would make the rule comparison a
comparison of travel models. The cost of the restriction is real ---
DAHS cannot exploit batching or routing synergies, and a controller that
jointly optimised release and batching could dominate it --- and we
record this in Section 8.

\subsection{Dispatching rules and data-centric control in
warehousing}\label{dispatching-rules-and-data-centric-control-in-warehousing}

Priority dispatching rules are the standard instrument for release
decisions under uncertainty: they are transparent, require no training,
and execute in the time available at a decision epoch. The classical
families and their properties are established (Conway et al. 1967;
Pinedo 2016) --- arrival-driven (FIFO), due-date-driven (EDD (Jackson
1955), minimum slack, modified due date (Baker and Bertrand 1982)),
processing-time-driven (WSPT (Smith 1956)), and composite
slack-and-processing indices (COVERT (Carroll 1965), ATC (Vepsäläinen
and Morton 1987)). The persistent finding across that literature is that
no single rule dominates: the rule minimising weighted tardiness under
light load is not the rule that does so once the queue saturates or the
due-date mix tightens (Pinedo 2016).

In warehousing specifically, order release and dispatching under
stochastic arrivals has been surveyed by Koster et al. (2007), Boysen et
al. (2019) and Boysen and Koster (2025), and the planning-problem
taxonomy of Gils et al. (2018) places dispatching among the operational
decisions that interact most strongly with due-date performance. The
field's recent direction is data-centric: Winkelhaus and Grosse (2020)
survey the shift toward data-driven warehouse control, and learned
controllers have been applied to dynamic order picking (Mahmoudinazlou
et al. 2025), order batching (Cheng et al. 2024), and job-shop
dispatching with an LSTM-PPO agent (Chen et al. 2024). What
distinguishes the warehouse setting from the job-shop setting these
methods were developed in is a second deadline clock on perishable
goods. We put that clock in the objective and \emph{measure} whether it
binds at a 15-minute review (Section 3.5). It does at the epoch level;
it is small at the order level (\(1.4\%\) of orders, \(1.3\%\) of
economic weight). The customer clock remains the dominant pivot.

\subsection{Simulation-trained selectors of dispatching
rules}\label{simulation-trained-selectors-of-dispatching-rules}

Selecting a rule as a function of the shop state, with the selector
trained on simulation output, is a mature line of work and predates this
paper by several decades. Wu and Wysk (1988) introduced the
\emph{multi-pass} construction: at a decision point, simulate the
candidate rules forward, record their outcomes, and use the result to
choose. Mouelhi-Chibani and Pierreval (2010) made the selector a learned
function, training a neural network on simulated states labelled by the
best-performing rule so that the multi-pass simulation is paid offline
and the deployed decision is a forward pass. Shiue et al. (2020) extend
the same construction with reinforcement learning over the rule set.
Durašević and Jakobović (2022) is a classifier over GP-evolved rules in
one unrelated-parallel-machines environment, not a survey. Surveys of
hyper-heuristics generally (Drake et al. 2020; Dokeroglu et al. 2024)
organise this literature around the offline-learning, online-application
paradigm that the present method also follows.

We correct a characterisation of this literature. Genetic programming in
this literature is used predominantly to \emph{generate} the low-level
rules that are subsequently selected among, not to learn the selector
(Branke et al. 2016; Nguyen et al. 2017); and prior selectors, including
those in Durašević and Jakobović (2022) and Mouelhi-Chibani and
Pierreval (2010), are trained essentially by supervised learning on
simulation-derived labels, which is what we do as well. Shiue et al.
(2020) is reinforcement learning over a rule set, not that supervised
construction. The method in this paper belongs inside that tradition
rather than departing from it.

\subsection{Rollout and classification-based approximate policy
iteration}\label{rollout-and-classification-based-approximate-policy-iteration}

A rollout policy improves a base policy by simulating each action at the
current state, following the base policy thereafter, and taking the
action of least simulated cost (Bertsekas 2020). Rollouts are a
cornerstone of approximate dynamic programming (Powell 2022) and are
typically truncated to a finite horizon for tractability (Bertsekas
2020; Goodson et al. 2017).

The construction we use (estimate action values by simulation at a
sample of states, then fit a classifier to represent the improved
policy) is the same family as \textbf{Rollout Classification Policy
Iteration}, introduced by Lagoudakis and Parr (2003) and developed by
Fern et al. (2006), Dimitrakakis and Lagoudakis (2008) and Farahmand et
al. (2015). It is not a new training paradigm. Offline rollout
generation for supervised policy learning has existed for over two
decades in the reinforcement-learning literature and, as Section 2.3
records, for longer than that in the scheduling literature. We also do
not run the policy-iteration loop: labels are generated once under a
behaviour policy and the ranker is fitted once, so the deployed
controller is a one-shot imitation of truncated rollouts rather than
iterated RCPI.

Two details of our instantiation differ from the standard RCPI setting,
and we note them as details rather than as contributions. The classifier
is fitted to the full per-action cost \emph{vector} rather than to the
arg-max alone, encoded as a tempered-softmax label distribution (Geng
2016); Section 6.8 reports an ablation showing this makes no material
difference, which is consistent with RCPI practice. And the rollout is
truncated at a short horizon, which is a standard truncation in the
rollout tradition (Bertsekas 2020).

\subsubsection{Rollout and ADP for dynamic
dispatching}\label{rollout-and-adp-for-dynamic-dispatching}

The operations-research literature on dynamic dispatching under
stochastic arrivals is the closest methodological neighbour to this
work. Klapp et al. (2018b) and Klapp et al. (2018a) formulate the
dispatch-waves problem --- when to release accumulated demand, given
that more will arrive --- a neighbouring dynamic-dispatching decision
(when to release a wave), not the inner rule-selection decision studied
here. Goodson et al. (2017) give a general rollout framework for
finite-horizon stochastic dynamic programs, including the treatment of
truncated horizons and pre- versus post-decision rollouts (Bertsekas
2020); Goodson et al. (2016) apply rollout policies to stochastic-demand
routing.

Ulmer et al. (2020) set out the modelling conventions for stochastic
dynamic routing and dispatching problems that are the closest
operational neighbour to Section 3. Section 3 itself uses Powell's
sequential-decision notation (Powell 2019), not route-based MDPs. The
same group's work on offline--online approximate dynamic programming
(Ulmer et al. 2019), on budgeting decision time under stochastic
requests (Ulmer et al. 2018), and on anticipation against reactive
re-optimisation (Ulmer 2019) is directly relevant to the trade-off this
paper studies.

One thread deserves particular emphasis because it supersedes part of
our construction. That literature does not simply truncate a rollout and
discard the tail: it \textbf{approximates the remainder with a learned
value function}, and in places uses learning to set the rollout horizon
state-dependently. Truncating the tail and discarding it is the worse
construction: a value-approximated remainder need not grow with the
remaining horizon, and would likely permit a shorter \(\tau\). We do not
implement it, and we record it in Section 9 as the most promising
extension.

\subsubsection{How this differs from value-function
approximation}\label{how-this-differs-from-value-function-approximation}

It is fair to ask what distinguishes this from value-function
approximation or reinforcement learning, since those also learn offline
from simulation. The method is \textbf{inside} the
approximate-dynamic-programming family, not outside it: what is
described here is one step of approximate policy iteration in which the
improved policy is represented by a classifier rather than derived from
a value function (Lagoudakis and Parr 2003; Fern et al. 2006).

\emph{What is learned.} Value-function approximation fits \(V(S)\) or
\(Q(S,u)\) --- a scalar satisfying, approximately, a Bellman fixed point
--- and recovers a policy from it by one-step lookahead. Here no value
function is ever formed and no fixed point is sought. The object fitted
is the policy itself: a classifier over actions, trained on directly
measured per-action costs at sampled states.

\emph{How error behaves.} This is the difference that matters and it is
why the comparison in Section 6.10 is informative. A bootstrapped value
estimate propagates error through the Bellman backup, so an error at one
state contaminates its predecessors and the fixed point can be reached
badly or not at all. The classification construction has no backup: its
error at each state is ordinary supervised generalisation error, and
errors at different states do not compound. That is the mechanism our
experiments are designed to isolate, holding the environment, corpus,
model class, feature set and objective fixed.

\emph{What each requires.} This asymmetry runs against us and we state
it plainly. Value-function approximation can be fitted from logged
transitions alone --- the data an operating warehouse already produces.
The construction here needs a simulator that can be \textbf{reset to an
arbitrary state and rolled forward under a counterfactual action},
because the label is the cost of rules that were not taken. That is a
strictly stronger requirement, it is the reason the method is
simulator-bound, and it is why the circularity of Section 8.1 is
intrinsic rather than incidental. A practitioner without a trustworthy
simulator should prefer value learning from logs; the comparison in this
paper is only relevant to one who has such a simulator and is deciding
how to use it.

\subsection{Reinforcement learning for
dispatching}\label{reinforcement-learning-for-dispatching}

Deep reinforcement learning has been applied widely to dispatching and
its scalability for production scheduling is under active study
(Stöckermann et al. 2025; Tassel et al. 2023). Two learning-based
selectors are the closest comparators and both appear as baselines here.
Imitation learning of dispatching decisions (Han and Jung 2025) trains
on the actions of a single expert dispatcher, and so requires an expert;
the multi-pass construction instead measures the counterfactual cost of
every rule and needs none. Offline reinforcement learning with maskable
action-value learning (Remmerden et al. 2025) learns a value function
from logged data, and is reimplemented faithfully in Section 5 as a
fitted-Q baseline (Ernst et al. 2005) trained on the same logged shifts.
We also include Proximal Policy Optimization (Schulman et al. 2017)
under an 8000-timestep budget (Stable-Baselines3
\texttt{total\_timesteps}, not review epochs), with a hyperparameter
sensitivity analysis in Section 6.9. A recent review (Liu et al. 2025)
frames simulation-derived self-labelling as an emerging paradigm for
machine learning in scheduling.

\subsection{Positioning}\label{positioning}

Given Sections 2.3 and 2.4, the mechanism at the centre of this paper is
not novel, and we do not claim it. Simulating a rule pool offline and
fitting a classifier to the result is RCPI in the reinforcement-learning
literature and multi-pass rule selection in the scheduling literature.
What we offer is an empirical study of truncated-rollout labels versus
online lookahead, fitted Q and PPO on one shift corpus, not a new
selector architecture.

Holding the environment, the seeds and the objective fixed, we vary how
the training signal is constructed: a directly measured per-action cost
vector, a bootstrapped state-action value fitted from the same logged
transitions on the same gradient-boosted-tree class, and a neural PPO
policy. PPO is not a matched simulation-step budget (Section 5). Section
6.10 reports the comparison; Section 6.9 reports the PPO hyperparameter
grid. The warehouse instantiation carries two deadline clocks; Section
3.5 measures whether the product clock binds at this review interval. On
this corpus the rollout-trained selector saturates by 50 training shifts
(Section 6.3). The matched FQI budget curve in Section 6.10 is the
pre-admit evaluation and is labelled as such.

The comparison is not a claim that the student recovers the teacher,
that extra rollout depth buys deployed \(J\), or that the selector is
robust under configuration transfer. The default scoreboard, the horizon
sweep and the robustness grid report the opposite on each of those
points. We make no claim to a new training paradigm.

\begin{center}\rule{0.5\linewidth}{0.5pt}\end{center}

\section{Problem Setting and
Simulator}\label{problem-setting-and-simulator}

\subsection{Orders, and the two deadline
clocks}\label{orders-and-the-two-deadline-clocks}

We consider a single warehouse shift of length \(T\) (8 hours). Orders
arrive over the shift. Order \(o\) carries

{\def\LTcaptype{none} % do not increment counter
\begin{longtable}[]{@{}
  >{\raggedright\arraybackslash}p{(\linewidth - 2\tabcolsep) * \real{0.5000}}
  >{\raggedright\arraybackslash}p{(\linewidth - 2\tabcolsep) * \real{0.5000}}@{}}
\toprule\noalign{}
\begin{minipage}[b]{\linewidth}\raggedright
symbol
\end{minipage} & \begin{minipage}[b]{\linewidth}\raggedright
meaning
\end{minipage} \\
\midrule\noalign{}
\endhead
\bottomrule\noalign{}
\endlastfoot
\(a_o\) & arrival time \\
\(p_o\) & processing time --- one pick tour, aggregating travel, search
and pick (Section 2.1) \\
\(d_o\) & \textbf{customer deadline}: the absolute time by which the
order must ship \\
\(x_o\) & \textbf{product deadline}: the absolute time after which the
goods are no longer saleable. Finite for perishable orders,
\(x_o = \infty\) otherwise \\
\(w_o\) & economic weight of the order, set by its priority class \\
\end{longtable}
}

The two deadlines are distinct constraints with distinct failure modes,
and they are sampled independently: missing \(d_o\) ships an order late,
missing \(x_o\) destroys the goods. Either can bind first. FEFO sorts on
\(x_o\); EDD sorts on \(d_o\). Section 3.5 tests whether the product
clock binds at this review interval.

\textbf{Where the expiry of an order comes from}. FEFO is conventionally
a rule for issuing inventory \emph{lots}, not customer orders, so an
order-level expiry requires justification. We model the stage
\emph{after} lot allocation: an upstream allocation policy --- exogenous
here, in the same sense as routing and batching (Section 2.1) --- has
already committed specific lots to specific orders, so each order
inherits a concrete expiry from the stock reserved against it. For a
single-line order that is the allocated lot's expiry. For a multi-line
order, the order ships as a unit and is therefore constrained by its
most perishable component, so \(x_o = \min_{\ell \in o} x_\ell\) over
the allocated lines \(\ell\). The experiments in this paper use
single-line orders, so \(x_o\) is the allocated lot's expiry directly;
the minimum-over-lines rule is the generalisation and requires no change
to the controller, since the controller reads only \(x_o\).

A fixed set of \(m = 10\) pickers processes orders; a picker handles one
order at a time and is busy for its processing time. Decisions are taken
at the boundaries of \(N = 32\) equal review intervals of \(L = 15\)
minutes. At each boundary the controller observes the system and chooses
a \emph{dispatching rule} from a fixed pool; that rule ranks the waiting
queue, and pickers are assigned down the ranking.

\subsection{The decision process, and what the controller can
see}\label{the-decision-process-and-what-the-controller-can-see}

We state the problem as a \textbf{sequential decision process} in the
canonical form of Powell (2019) and Powell (2022): state, decision,
exogenous information, transition function, objective, before any
implementation detail, and Section 3.4 then specifies the transition
concretely. Readers who prefer the reinforcement-learning vocabulary may
read the five elements as a finite-horizon Markov decision process; the
partial-observability caveat below applies under either reading.

\textbf{The sequential decision process}. Decisions are taken at epochs
\(t = 0, 1, \dots, N-1\), one per review interval. The \textbf{state} is

\[ S_t \;=\; \big(\, \mathcal{Q}_t,\; \mathbf{b}_t,\; t \,\big), \]

where \(\mathcal{Q}_t\) is the multiset of waiting orders, \emph{each
carrying its full attribute tuple} \((a_o, p_o, d_o, x_o, w_o)\), and
\(\mathbf{b}_t \in \mathbb{R}^m\) records when each picker next becomes
free. The \textbf{decision} is \(u_t \in \mathcal{H}\), a rule from the
pool. The \textbf{exogenous information} \(W_{t+1}\) is the set of
orders arriving in \((t, t+1]\) with their attributes. The
\textbf{transition} \(S_{t+1} = S^M(S_t, u_t, W_{t+1})\) is the
admission--rank--assign procedure of Section 3.4. Writing
\(C(S_t, u_t, W_{t+1})\) for the cost accrued over the interval (Section
3.3), the problem is

\[ \min_{\pi \in \Pi} \; \mathbb{E}\left[\; \sum_{t=0}^{N-1} C\big(S_t,\, U^\pi_t(S_t),\, W_{t+1}\big) \;\right] . \]

\textbf{The observation is not the state}. The controller does not see
\(S_t\). It sees

\[ x_t \;=\; \phi(S_t) \;\in\; \mathbb{R}^{26}, \]

a fixed-length summary listed in Appendix A. Calling \(x_t\) ``the
state'' is wrong as terminology and, more importantly, as mechanism:
\(\phi\) records \emph{marginal} summaries --- queue length, mean and
standard deviation of slack, mean processing time, counts of critical
and perishable orders --- and discards the \emph{joint} distribution
over per-order attributes. But it is the joint distribution that
determines what a ranking rule does next, because a rule orders orders
by a function of their attributes taken together.

\textbf{A witness}. \(\phi\) is not injective. The constructed pair
below is a searched-grid example, not a claim that every coordinate is
attained under the deployed processing-time law. Consider two queues of
two orders each, arriving at the same instant, differing only in which
order carries the tight deadline:

\[ \mathcal{Q}^{A} = \{(p_{\text{short}}, s_{\text{loose}}),\, (p_{\text{long}}, s_{\text{tight}})\}, \qquad \mathcal{Q}^{B} = \{(p_{\text{short}}, s_{\text{tight}}),\, (p_{\text{long}}, s_{\text{loose}})\}, \]

writing \(s_o = d_o - t - p_o\) for slack. The two have the same queue
length, the same arrival times and therefore the same ages, the same
mean and standard deviation of slack, the same mean processing time, and
--- with \(s\) and \(p\) chosen so that the critical thresholds fall the
same way --- the same critical and at-risk counts. Every coordinate of
\(\phi\) agrees: \(\phi(S^A) = \phi(S^B)\) exactly. The dynamics do not
agree. In \(\mathcal{Q}^{A}\) the binding deadline sits on the order
that occupies a picker longest, so deferring it is expensive and the
feasible set of on-time completions is strictly smaller than in
\(\mathcal{Q}^{B}\).

Two conditions are needed for that difference to be \emph{realised}, and
both are part of the construction rather than incidental to it.

\emph{The queue must contend for a picker.} A ranking expresses a
preference only when something has to wait. With the deployed ten
pickers and a two-order queue both orders start immediately whatever the
ranking says, every rule produces the identical trajectory, and the cost
gap is zero by construction. The witness is therefore built at one
picker. This is not a weakening: contention is the regime in which rule
selection has any effect at all, so it is the only regime in which
partial observability can cost anything.

\emph{The rule must key on slack and processing time jointly.} Queues
\(\mathcal{Q}^A\) and \(\mathcal{Q}^B\) carry the \emph{same multiset}
of slacks, so a rule that ranks only on a slack-like clock will often
fail to separate them. EDD sorts by due date, not slack; EEDD, MS and
MDD are the closer slack-family rules. No witness is claimed for those
rules. The composite rules ATC and COVERT rank on slack \emph{and}
\(p_o\) together, which is precisely the interaction \(\phi\) discards
by recording the two marginals separately. That is the sharper statement
of the defect: \(\phi\) retains the marginal distributions of slack and
of processing time and destroys their coupling.

\texttt{experiments/observability\_analysis.py} searches a grid of
picker counts, processing times and slacks over every rule in the pool,
verifies that the feature vectors coincide to machine precision before
comparing anything, and reports the strongest gap. The largest ATC gap
on that grid is \(3.79\) against \(-0.01\) at
\((p_{\mathrm{short}}, p_{\mathrm{long}}) = (4, 18)\) with one picker.
Processing time in the simulator is Triangular\((2,5,12)\), so \(p=18\)
is off-support. The only on-support ATC cell with a nonzero gap is
\((2,12)\), slacks \((-5,30)\), gap \(0.40\) (already-late tight job).
Displayed slacks \((0,40)\) have gap \(0\) on-support. The pair shows
that \(\phi\) can match while ATC cost depends on the latent pairing. It
is not a reachable-state regret floor, and it does not exhibit
conflicting arg-mins. \(\phi\) is still not a sufficient statistic.

\textbf{Consequence: a POMDP, and a policy-function approximation}. We
therefore do not claim that \(\phi\) is a sufficient statistic --- it is
not, and the witness settles it. The control problem is a
\textbf{partially observed} Markov decision process, and the policy
class we search is a policy-function approximation over the observation,

\[ U^\pi(S_t) \;=\; \arg\max_{h \in \mathcal{H}} \; f_\theta\big(\phi(S_t)\big)_h , \]

with no belief state maintained and no history beyond the lags \(\phi\)
carries explicitly. Two consequences follow, and they pull in different
directions.

The unfavourable one is that partial observability \emph{can} leave an
irreducible regret floor: two states that \(\phi\) cannot separate must
receive the same action, so whenever their optimal actions differ, some
regret is incurred that no amount of data or model capacity can remove.
The constructed pair does not exhibit that conflict. We treat the floor
as a possible consequence of \(\phi\), not as a result of the
off-support witness. Over the training corpus we locate mutual
near-neighbours in standardised \(\phi\)-space and report how often they
disagree about the cost-minimising rule, and what acting on the
neighbour's choice costs, as a share of the total benefit available from
rule selection. That number is an upper bound on the part of the
residual regret attributable to partial observability (\(3.7\%\) in
Section 8), and it is reported alongside the other limitations rather
than buried.

The favourable one is more specific than it first appears. The rollout
\textbf{labels} are computed from the true state \(S_t\): the simulator
holds \(\mathcal{Q}_t\) in full, and the rollouts of Section 4.3 run the
actual dynamics from it. So the supervision target is \emph{correct}; it
is only the \emph{covariates} that are lossy. The learning problem is
therefore regression with insufficient covariates --- a Bayes risk
induced by the feature map --- and not a misspecified or biased target.
This distinction matters for interpreting the results: a residual gap
between DAHS and a controller with access to \(S_t\) is attributable to
\(\phi\), not to the training signal, and enlarging \(\phi\) (a set
encoding over queued orders, for instance) is the remedy. We leave that
to future work and record it in Section 9.

\subsection{The objective}\label{the-objective}

Let \(f_o\) denote the completion time of order \(o\) if it is
dispatched. For an order still waiting at the reference horizon \(T\) we
set \(f_o = T + p_o\): the earliest it could possibly finish, since it
still requires a full pick. With that convention every order that
arrived has a well-defined outcome whether or not it was ever served,
and

\[ J \;=\; \sum_{o \in \mathcal{A}} w_o \Big[\, W_{b}\,\mathbb{1}\{f_o > d_o\} \;+\; W_{t}\,\max(f_o - d_o,\,0) \;+\; W_{s}\,\mathbb{1}\{o\in\mathcal{P}\}\,\mathbb{1}\{f_o > x_o\} \,\Big] \;+\; W_{h}\,|Q_T| , \]

where \(\mathcal{A}\) is the set of orders that arrived during the
shift, \(\mathcal{P}\subseteq\mathcal{A}\) are the perishable ones, and
\(|Q_T|\) is the queue length at shift end. Non-perishable orders have
no \(x_o\) and never attract \(W_s\). Weights are \(W_{b} = 3.0\) (late
shipment), \(W_{t} = 0.2\) per minute of lateness, \(W_{s} = 5.0\)
(spoilage), and \(W_{h} = 0.005\) per queued order. They are fixed
before any learning and are not tuned to any method.

Three features of this objective are consequential.

\textbf{Priority class now enters the objective}. WSPT and ATC rank by
\(w_o/p_o\). An unweighted objective would grade those rules against a
criterion the evaluation never measured (Section 6.2). Rule and
objective agree.

\textbf{Perishability now enters the objective}. The \(W_s\) term is the
only place a product deadline can be priced. Without it,
``perishability-constrained'' would not be a property of the
optimisation problem at all.

\textbf{Orders that are never served are charged}. Charging an abandoned
order only \(W_{h} = 0.005\) against \(W_{b} = 3.0\) for one served late
would give a controller a factor-of-600 incentive to decline difficult
orders, and a completed-orders-only breach rate would drop those orders
from the metric entirely. The convention \(f_o = T + p_o\) closes both
gaps: an unserved order past its deadline is charged exactly as a late
one, and the \(+\,p_o\) ensures that dispatching an order onto a free
picker costs precisely what abandoning it costs, with any earlier
dispatch costing strictly less. Doing the work is therefore weakly
optimal by construction. \(W_h\) survives strictly as a work-in-progress
holding cost.

\textbf{Spoilage mechanics, stated explicitly}.

\emph{Does a perishable order have a distinct expiry, or is it the due
date?} Distinct. \(x_o\) is drawn independently of \(d_o\) (Section
3.1), so for a perishable order either clock can bind first.

\emph{What happens when an order spoils?} Its goods become unsaleable at
\(x_o\). The simulator is discrete-interval: spoilage is assessed when
the potential or a KPI is computed (at each review and at the shift-end
reference \(T\)), not as a continuous-time event at \(x_o\). Objective
spoilage uses \(f_o > x_o\), with \(f_o = T + p_o\) if the order is
still waiting. Dispatch after expiry does not clear \(W_s\). The KPI
predicate \texttt{is\_expired\_at} is stricter for unserved orders: it
fires only when \(T > x_o\), with no \(+p_o\). The two can disagree ---
an unserved perishable with \(p_o=6\), \(x_o=105\) and \(T=100\) is
spoiled under \(J\) and not expired under the KPI. The order is
\textbf{not} removed from the queue: spoiled stock still has to be
pulled and disposed of, so it continues to consume a picker when it is
eventually handled. Keeping it in the queue also closes an incentive
gap. If spoiled orders vanished, a controller could free picking
capacity by stalling until perishables expired, which is the same class
of loophole as exempting unfinished orders.

\emph{Is a spoiled order counted in the breach count?} Lateness and
spoilage are \textbf{separate predicates} on the same order, and an
order may be neither, either, or both. When both fire, the composite
cost charges both --- they are distinct economic losses, a late shipment
and destroyed stock. In the reported metrics they are kept apart so
nothing is double-counted in a headline: the breach rate counts lateness
only, the spoilage rate counts spoilage only, and the primary
\textbf{service-failure rate} counts an order once if it is late
\emph{or} spoiled.

\emph{How are unfinished perishables penalised at shift end?} Through
the same \(f_o = T + p_o\) convention as any other unserved order
(above). If \(T + p_o > x_o\) the order is spoiled and charged
\(W_s w_o\); if \(T + p_o > d_o\) it is also late and charged
\(W_b w_o\) plus tardiness. Since shelf life is at most 120 minutes
against a 480-minute shift, a perishable order still waiting at shift
end is spoiled with near certainty, which is the intended behaviour:
abandoning perishable stock is the most expensive thing this objective
can do.

\textbf{How the reported rates are calculated}. Writing \(\mathcal{A}\)
for the set of orders admitted by the last review epoch (Section 3.4),
\(\mathcal{S} \subseteq \mathcal{A}\) for those dispatched, and
\(\mathcal{P} \subseteq \mathcal{A}\) for the perishable ones. The
composite cost uses \(f_o = T + p_o\) for any order still waiting at the
reference horizon \(T\). The reported \emph{rates} do not: an unserved
order is late or spoiled in the KPI only if the relevant clock has
already passed at \(T\), with no \(+p_o\). That split is deliberate.
Adding \(p_o\) to the KPI would count an unserved order whose due date
is still in the future as a service failure.

\[ \text{service-failure rate} = \frac{|\{o \in \mathcal{A} : \text{overdue or expired as of } T\}|}{|\mathcal{A}|}, \qquad \text{spoilage rate} = \frac{|\{o \in \mathcal{P} : \text{expired as of } T\}|}{|\mathcal{P}|}, \]

\[ \text{breach rate}_{\text{arrived}} = \frac{|\{o \in \mathcal{A} : \text{overdue as of } T\}|}{|\mathcal{A}|}, \qquad \text{breach rate}_{\text{served}} = \frac{|\{o \in \mathcal{S} : \text{overdue as of } T\}|}{|\mathcal{S}|} . \]

Served orders are scored at finish time; unserved orders are scored at
the clock \(T\), not at \(T+p_o\). ``As of \(T\)'' means that split.

The last of these is the completed-orders-only breach rate, reported
here under that explicit name. Its denominator excludes every order the
controller declined to dispatch, and with \(W_h = 0.005\) against
\(W_b = 3.0\) so did an objective that charged abandoned orders only the
holding term. Both rates are reported, under names that make the
denominator explicit. \(\mathcal{A}\) includes orders rejected at the
door when the queue was at capacity; they are real demand that went
unmet, and excluding them would reopen the same gap in a different
place. \(\mathcal{A}\) \textbf{does} include arrivals in the last open
interval \((T-L, T]\): admission runs at review epochs, and a terminal
admit at \(T\) enters every order with arrival time \(\le T\) into
\(\mathcal{A}\) as unserved. They are counted, never dispatched. Mean
\(|\mathcal{A}|=791\) against a Poisson mean of \(1.65\times 480=792\).
Those late arrivals are shared across every method.

\subsection{The simulator and its
parameters}\label{the-simulator-and-its-parameters}

The transition function \(S^M\) of Section 3.2 is realised as a
discrete-interval simulation model; its construction, verification and
validation follow standard practice (Law and Kelton 2000; Sargent 2013),
with the input distributions fitted to a public order trace where one
exists (Section 6.7) and the remaining parameters carrying explicit
provenance tags below.

The environment is a periodic-review discrete-interval model. Within
interval \([t, t+L)\) it (i) admits every order that has arrived
\textbf{by \(t\)}, subject to a queue capacity of 200, recording
overflow as \emph{rejected demand} rather than discarding it; (ii) ranks
the waiting queue with the chosen rule, evaluated at \(t\); (iii)
assigns down the ranking to the earliest-free picker with start time
\(\max(\text{picker free}, t)\), stopping once no picker can start
before \(t+L\); and (iv) records end-of-interval statistics. Orders
arriving inside the interval wait for the next epoch.

The admission rule matters. Admitting every order arriving before
\(t+L\) (fifteen minutes of look-ahead) and setting the start time to
\(\max(\text{picker free}, a_o, t)\) would reserve a picker for a
not-yet-arrived order and leave it idle until that order appeared. Rules
sorted by arrival never paid this cost; arrival-agnostic rules would pay
it constantly (Section 6.2).

\textbf{Parameters and their provenance}. Every input is now either
fitted to data, grounded in a cited source, or declared a design choice;
none is left unexplained.

{\def\LTcaptype{none} % do not increment counter
\begin{longtable}[]{@{}
  >{\raggedright\arraybackslash}p{(\linewidth - 4\tabcolsep) * \real{0.3333}}
  >{\raggedright\arraybackslash}p{(\linewidth - 4\tabcolsep) * \real{0.3333}}
  >{\raggedright\arraybackslash}p{(\linewidth - 4\tabcolsep) * \real{0.3333}}@{}}
\toprule\noalign{}
\begin{minipage}[b]{\linewidth}\raggedright
Input
\end{minipage} & \begin{minipage}[b]{\linewidth}\raggedright
Value
\end{minipage} & \begin{minipage}[b]{\linewidth}\raggedright
Provenance
\end{minipage} \\
\midrule\noalign{}
\endhead
\bottomrule\noalign{}
\endlastfoot
Shift, review interval, pickers & 8 h, 15 min, 10 & Operating point \\
Queue capacity & 200 orders & Operating point \\
Inter-arrival \textbf{shape} & Poisson (exponential inter-arrivals) &
Operating point. Olist shape is lognormal; Section 6.7 and Appendix C
report the fit and a frozen-ranker replay under empirical bootstrap
arrivals \\
Arrival \textbf{rate} & 1.65 orders/min nominal & Operating point; swept
in Section 6.5 \\
Processing time \(p_o\) & Triangular\((2, 5, 12)\) min &
\textbf{Literature}: three-point time standard for manual
picker-to-parts picking (Tompkins et al. 2010; Koster et al. 2007) \\
Customer window \(d_o - a_o\) & Triangular\((15, 45, 90)\) min &
\textbf{Operating point.} AIC on the Olist
purchase-to-estimated-delivery sample selects lognormal, not triangular
(Appendix C). The campaign uses this triangular envelope at warehouse
scale; it is not the MLE \\
Shelf life \(x_o - a_o\) & Triangular\((20, 60, 120)\) min & Design
parameter; no public trace carries expiry. Swept in Section 6.11 \\
Perishable fraction & 0.20 & Design parameter; varied by scenario \\
Priority classes and weights & \(\{\)low, medium, high\(\}\) at
\((0.50, 0.35, 0.15)\), \(w_o \in \{1, 2, 4\}\) & Design parameter \\
\end{longtable}
}

Fitting is the right operation on a trace that exists, and Section 6.7
reports the fits, the candidate families compared by AIC, and, for
processing time, where the trace carries no warehouse pick-time field,
the reason no fit is attempted.

\subsection{Does perishability bind at a 15-minute
horizon?}\label{does-perishability-bind-at-a-15-minute-horizon}

A model can carry a product deadline without that deadline ever changing
a decision. Since we claim the setting is perishability-constrained, we
test the claim directly rather than assert it. The criterion is
decision-relevance: in what fraction of decisions does delaying an order
by one review interval alter its outcome?

At each epoch \(t\), a waiting order has exactly two options available
to the dispatcher --- served now, completing at \(t + p_o\), or
deferred, completing no earlier than \(t + L + p_o\). Call the order
\textbf{expiry-pivotal} at \(t\) when those two straddle its product
deadline,

\[ t + p_o \;\le\; x_o \;<\; t + L + p_o , \]

so that one interval of delay is the difference between saleable goods
and waste, and \textbf{due-pivotal} when the same holds for \(d_o\).

Three conditions were fixed \emph{before} the diagnostic was run, and
all three had to hold for the framing to stand: at least 5\% of
decisions carrying an expiry-pivotal order; at least 10\% of perishables
having their product clock bind before their customer clock; and the
choice of rule moving the spoilage count at at least 10\% of epochs. Had
any failed, the correct response was to drop the perishability framing
and the expiry-aware rules, and to report that instead.

\textbf{Table 1}. Perishability decision-relevance, 30 calibration
shifts × 8 behaviour rules (the nine-candidate \texttt{SCREENING\_POOL}
without FIFO: EDD, EEDD, FEFO, WSPT, ATC, MS, MDD, COVERT). \(7{,}440\)
recorded decision epochs (empty queues omitted). These are repeated
measures on the same 30 shifts, not \(7{,}440\) independent shifts. The
shipped deployed pool \texttt{resolve\_pool} is six rules; reproducing
this table from current \texttt{config.yaml} will not recover it.

{\def\LTcaptype{none} % do not increment counter
\begin{longtable}[]{@{}
  >{\raggedright\arraybackslash}p{(\linewidth - 4\tabcolsep) * \real{0.2727}}
  >{\raggedleft\arraybackslash}p{(\linewidth - 4\tabcolsep) * \real{0.3636}}
  >{\raggedleft\arraybackslash}p{(\linewidth - 4\tabcolsep) * \real{0.3636}}@{}}
\toprule\noalign{}
\begin{minipage}[b]{\linewidth}\raggedright
Quantity
\end{minipage} & \begin{minipage}[b]{\linewidth}\raggedleft
Measured
\end{minipage} & \begin{minipage}[b]{\linewidth}\raggedleft
Threshold
\end{minipage} \\
\midrule\noalign{}
\endhead
\bottomrule\noalign{}
\endlastfoot
Decisions with an expiry-pivotal order in queue & \textbf{35.5\%} & ≥
5\% \\
Perishables whose expiry binds before their due date & \textbf{27.6\%} &
≥ 10\% \\
Epochs where the rule choice changes the spoilage count &
\textbf{91.0\%} & ≥ 10\% \\
Queued perishables that are expiry-pivotal & 7.3\% & --- \\
All orders that are expiry-pivotal & 1.4\% & --- \\
Share of economic weight on expiry-pivotal orders & 1.3\% & --- \\
Mean spoilage \emph{weight} spread across rules, when discriminating &
3.97 & --- \\
\end{longtable}
}

All three conditions are met at perishable fraction \(0.20\), so the
constraint is real at this horizon and the framing stands at the
operating point we ran. We did not rerun Table 1 at \(0.05\) or at
\(0.01\). The Olist food/drink share is \(0.0099\), about one twentieth
of \(0.20\). If the 35.5\% expiry-pivotal rate scaled linearly with the
perishable fraction, it would sit near \(1.8\%\) at \(0.01\) and would
miss the 5\% mark. That scaling is untested. Section 6.11 sweeps shelf
life by \(\times\{0.8,1.0,1.25\}\) around this inflated fraction; it is
not a perishable-fraction sweep.

Two qualifications belong with that conclusion rather than after it. The
\emph{marginal} rate is small --- only 1.4\% of individual orders are
expiry-pivotal at any given epoch, carrying 1.3\% of economic weight ---
so perishability is not the dominant cost driver; the customer clock is,
with 95.1\% of epochs carrying a due-pivotal order against 35.5\%
carrying an expiry-pivotal one. The \(3.97\) figure is
\texttt{mean\_spoilage\_spread\_when\_discriminating} from
\texttt{pivotality\_summary.json}: a mean difference of economic
\emph{weight} (\(w_o\in\{1,2,4\}\)), not an order count. The diagnostic
stores the weight spread and discards the count. What makes the clock
decision-relevant is concentration and frequency: the pivotal orders
cluster, so a third of all decisions have at least one in the queue, and
at 91.0\% of epochs the choice of rule moves realised spoilage weight. A
constraint that changes the outcome at nine epochs in ten is binding on
the controller whatever share of orders it touches.

\subsection{The rule pool}\label{the-rule-pool}

The pool is not a convenience sample. Candidates span the four
information sources a dispatcher can key on, so that ``no single rule
dominates'' is a structural property of the design space rather than an
empirical accident. Nine candidates are screened (Section 6.1); the six
that survive are marked.

{\def\LTcaptype{none} % do not increment counter
\begin{longtable}[]{@{}
  >{\raggedright\arraybackslash}p{(\linewidth - 6\tabcolsep) * \real{0.2308}}
  >{\raggedright\arraybackslash}p{(\linewidth - 6\tabcolsep) * \real{0.2308}}
  >{\raggedright\arraybackslash}p{(\linewidth - 6\tabcolsep) * \real{0.2308}}
  >{\centering\arraybackslash}p{(\linewidth - 6\tabcolsep) * \real{0.3077}}@{}}
\toprule\noalign{}
\begin{minipage}[b]{\linewidth}\raggedright
Rule
\end{minipage} & \begin{minipage}[b]{\linewidth}\raggedright
Keys on
\end{minipage} & \begin{minipage}[b]{\linewidth}\raggedright
Source
\end{minipage} & \begin{minipage}[b]{\linewidth}\centering
Retained
\end{minipage} \\
\midrule\noalign{}
\endhead
\bottomrule\noalign{}
\endlastfoot
FIFO & arrival only & zero-information control & --- \\
EDD & customer deadline & Jackson (1955) & ✓ \\
\textbf{EEDD} & \textbf{both deadlines}: \(\min(d_o, x_o)\) & this paper
(see below) & ✓ \\
MS & customer deadline slack & Conway et al. (1967) & ✓ \\
MDD & customer deadline, degrading to SPT when past due & Baker and
Bertrand (1982) & ✓ \\
FEFO & product deadline only & classical lot-issuing rule & --- \\
WSPT & weight and processing time & Smith (1956) & --- \\
ATC & slack × processing, exponential discount & Vepsäläinen and Morton
(1987) & ✓ \\
COVERT & slack × processing, linear truncation & Carroll (1965) & ✓ \\
\end{longtable}
}

Four points about the pool.

\textbf{FEFO is not a due-date rule, and neither rule alone is right}.
FEFO sorts on \(x_o\); EDD sorts on \(d_o\). That is EDD's clock, not
FEFO's. With two deadlines, \emph{neither} single-clock rule is the
sensible one. EDD ignores expiry entirely. FEFO ranks on \(x_o\), which
is \(\infty\) for the 80\% of orders that are not perishable, so it
dumps every non-perishable order behind every perishable one --- on this
order mix that is close to a strawman, and it is why the FEFO mask
(Section 4.3) had to exist at all.

The rule a scheduler would actually write, told that orders carry two
deadlines, sorts ascending on the \textbf{effective deadline}
\(\min(d_o, x_o)\) --- whichever clock binds first. We call it
\textbf{EEDD}, and it is in the pool because leaving it out would have
made ``expiry-awareness buys nothing'' an artefact of testing only a
degenerate expiry rule. Section 6.1 reports what it earns.

\textbf{ATC is calibrated, not assumed}. WSPT is exactly the
\(k \to \infty\) limit of ATC: as the look-ahead scale grows,
\(\exp(-\text{slack}/(k\bar p)) \to 1\), leaving the WSPT index
\(w_o/p_o\). A correctly calibrated ATC therefore cannot be beaten by
WSPT, and an unfitted \(k\) fixed at 2.0 would invert that ordering. We
calibrate \(k\) twice, on a calibration corpus disjoint from both
training and test shifts: once for \textbf{standalone} use, because ATC
is itself a reported benchmark and benchmarking an uncalibrated rule
understates it; and once for \textbf{portfolio contribution}, which is
the quantity that matters when the rule sits inside a selector. COVERT's
scale is calibrated the same way. Section 6.1 reports both values and
the deployed one.

\textbf{Screening is by marginal contribution, not win rate}. A rule
earns its slot by covering states the others handle badly. We therefore
report, per rule, both its win rate and the increase in achievable cost
if it were removed from the pool, with a bootstrap interval on the
latter. A rule with a high win rate and zero marginal contribution is
redundant; a rule with a low win rate and positive marginal contribution
is a specialist worth keeping. Win rate alone cannot distinguish the
two. This is also how we answer what \textbf{FIFO} contributes in a
due-date-driven setting: it enters as the zero-information control and
the screen reports whether it earns its place.

\begin{center}\rule{0.5\linewidth}{0.5pt}\end{center}

\section{The DAHS Method}\label{the-dahs-method}

\subsection{Overview}\label{overview}

DAHS is an offline-learned, online-applied selection hyper-heuristic.
Training proceeds in four stages: (1) generate a corpus of simulated
shifts under a state-covering behaviour policy; (2) for every decision
state, run a truncated-horizon rollout of each rule and record the
per-rule cost vector, from which a training label is formed; (3)
discover a small set of operating \emph{regimes} and append
regime-membership features; (4) fit a calibrated gradient-boosted ranker
to the rollout-derived labels. At deployment a lightweight
\emph{switching controller} wraps the ranker with a dwell constraint and
an entropy gate. The expiry-rule mask is a no-op on the screened pool
(Section 4.3). We describe each stage in turn.

\subsection{The observation vector}\label{the-observation-vector}

At each decision epoch the controller receives \(\phi(S_t)\), a
26-dimensional summary of the true state (Section 3.2) covering queue,
resource, customer- deadline, product-deadline, arrival, history and
temporal context. It is an observation, not a state, and Section 3.2
gives the explicit pair of queues it cannot separate. Appendix A lists
every feature with its group and the source or design rationale.

Two features that would make the matrix singular are omitted.
\texttt{time\_to\_next\_expected\_carrier} is \(1/\lambda\) and
therefore \textbf{constant} within a configuration;
\texttt{intervals\_remaining} is an exact affine function of
\texttt{interval\_index\_in\_shift}, the two summing to \(N\) by
construction. Together they make the feature matrix exactly singular,
which silently corrupts the regime layer's model selection (Section
4.4). Three expiry-pressure features are included, since the product
deadline enters the objective and a selector cannot act on a constraint
it cannot see. Appendix A also reports the correlation and
variance-inflation analysis, and Section 6.8 records a test-set SHAP
diagnostic that is not treated as a held-out ablation.

To \(\phi(S_t)\) DAHS appends the regime-membership posteriors of
Section 4.4.

\subsection{Rollout-informed training
labels}\label{rollout-informed-training-labels}

The supervisory signal is generated as follows. We simulate a corpus of
training shifts under a state-covering behaviour policy, giving one
decision state per review epoch per shift. For each state \(s_t\) we
form a label over the pool by \textbf{multi-sample truncated rollout}:

\begin{enumerate}
\def\labelenumi{\arabic{enumi}.}
\tightlist
\item
  Walk the shift forward to epoch \(t\), so that \(s_t\) is the true
  state \(S_t\) with its full queue.
\item
  Draw \(M\) independent continuations. Continuation \(m\) freezes the
  realised history at \(t\) and resamples the \emph{unrealised} future
  --- arrivals after \(t\) and their attributes --- from a stream seeded
  by \texttt{SeedSequence({[}base\_seed,\ shift\_seed,\ t,\ m{]})}, not
  on the rule under test.
\item
  For each rule \(h\), hold \(h\) fixed for the next \(\tau\) epochs of
  continuation \(m\) (open-loop truncated simulation of that rule ---
  not one improving action followed by a different base policy), and
  record the cost \(\hat{J}^{\tau}_{h,m}(s_t)\) accrued over that
  window. Average:
  \[ \hat{J}^{\tau}_h(s_t) \;=\; \frac{1}{M}\sum_{m=1}^{M} \hat{J}^{\tau}_{h,m}(s_t), \qquad \widehat{\mathrm{se}}_h(s_t) \;=\; \frac{\hat{\sigma}_h(s_t)}{\sqrt{M}} . \]
\item
  Convert the cost vector into a probability distribution by a
  \textbf{tempered softmax} at a state-dependent temperature:
  \[ p^{\tau}_h(s_t) = \frac{\exp(-\hat{J}^{\tau}_h(s_t)/T(s_t))}{\sum_{h'} \exp(-\hat{J}^{\tau}_{h'}(s_t)/T(s_t))}, \qquad T(s_t) = \beta\,\hat\sigma(s_t), \]
  where \(\hat\sigma(s_t)\) is the standard deviation of that state's
  own cost vector across the pool and \(\beta\) is a single
  dimensionless multiplier fitted once on the training corpus.
\end{enumerate}

\textbf{Why \(M > 1\)}. If every stochastic quantity is pre-sampled when
a shift is generated, every rule sees the identical realised future
belonging to the shift seed. The label then records which rule was best
\textbf{in hindsight on one path}, not which had the lowest expected
cost, and the rollout variance is identically zero. Hindsight-optimal on
one path and lowest-in-expectation are different quantities. The
estimator below records a per-cell standard error
\(\widehat{\mathrm{se}}_h\) alongside every label so the residual noise
is reported rather than assumed away.

\textbf{Terminal admit and the labels.} Live evaluation admits leftover
arrivals at \(T\) (\texttt{sim.terminal\_admit:\ true}). Production
labels were generated with that flag off and were not regenerated for
the completeness run. Truncated windows that exhaust the remaining shift
\emph{do} call the terminal admit when the flag is on, so the guarantee
that truncated rollouts are independent of the flag is false for
last-\(\tau\) epochs. The default comparison in Section 6.2 is an
eval-only refresh of frozen controllers, not a relabelled campaign.

\textbf{Common random numbers}. The continuation seed depends on the
shift, the epoch and the sample index, and deliberately \textbf{not} on
the rule under test. Every candidate is therefore scored against an
identical set of \(M\) futures. The comparison is paired, and since the
tempered softmax reads only \emph{differences} between rules, the
relevant variance is that of \(\hat{J}^{\tau}_h - \hat{J}^{\tau}_{h'}\),
which is far smaller than that of either term: the shared arrival shock
that dominates a single rollout's cost cancels. This is why a modest
\(M\) suffices where independent sampling would need an order of
magnitude more. Section 6.4 sweeps \(M\) and reports where the label
stabilises.

\textbf{Cost}. Walking each shift forward once and branching at each
epoch costs \(O(N \cdot |\mathcal{H}| \cdot M \cdot \tau)\)
interval-steps per shift. A labeller that replayed each shift from
\(t = 0\) for every epoch and every rule would be quadratic in \(N\).
The linear walk is what makes \(M>1\) affordable. Section 6.12 reports
the measured totals.

\textbf{Why the temperature is per state}. A single temperature fitted
across the whole corpus is dominated by expensive late states. Charging
orders that are never served (Section 3.3) makes a state's cost scale
with its outstanding work, so the spread of the per-rule cost vector
grows through a shift --- by two orders of magnitude between the first
epochs and the last, measured on the corpus. A single temperature is
then set by the expensive late states, and the early ones, where the
rules genuinely differ by little, come out close to uniform; the median
label entropy sits above any target band at every multiplier. Dividing
by the state's own spread asks the question the ranker actually has to
answer --- \emph{how much better is the best rule than the others, here}
--- and leaves the absolute cost of the state, which the ranker does not
predict, out of the label. States where every rule ties carry no
information at any temperature and fall back to the uniform.

\(\beta\) is then selected once by a one-dimensional bisection on the
training corpus, so that the median training-label entropy falls in a
target band expressed as a fraction of \(\log|\mathcal{H}|\) rather than
in absolute nats --- sharp enough to be informative, soft enough to
retain the cost margin. On the regenerated corpus (\(|\mathcal{H}|=6\),
\(M=20\), \(\tau=4\), 250 train shifts) the search selected
\(\beta = 0.470\) under per-row temperature and achieved median
train-row entropy \(0.638\) nats against the band \([0.387, 0.905]\).
The reported value is not a node of \texttt{beta\_grid}. The same
\(\beta\) is applied to the test and calibration corpora, which would
otherwise not be on the ranker's scale.

Two corrections are applied consistently in both labelling and
deployment: when the perishable fraction is below 0.05 the FEFO mass is
zeroed and the distribution renormalised, because FEFO cannot rank a
queue with no product deadlines; and, \emph{for the test corpus only},
states whose maximum label probability falls below
\(\theta = 2.2/|\mathcal{H}|\) are filtered out as genuinely ambiguous
decisions. The threshold scales with the pool for the same reason the
entropy band does.

\textbf{The mask is inert in the deployed model, and we leave it in
place rather than delete it.} Screening dropped FEFO (Section 6.1), so
there is no expiry-only rule left for the mask to suppress and it is a
no-op on the deployed pool. It is retained because it is a property of
the \emph{labelling and deployment contract}, not of one pool: any pool
containing a rule that ranks on \(x_o\) alone needs it, and removing the
code would silently reintroduce the defect if FEFO ever returned. It
should not, however, be counted as a working component of the method as
deployed. On the screened pool it is a no-op. EEDD, which ranks on
\(\min(d_o, x_o)\), needs no mask: on a queue with no perishables it
degrades continuously to EDD rather than becoming undefined.

\textbf{Where the corpora come from}. Shift seeds are drawn once from a
single \texttt{SeedSequence} and partitioned into three contiguous,
disjoint blocks: training, calibration and test. The calibration block
exists so that rule hyperparameters (ATC's and COVERT's look-ahead
scales, Section 3.6) can be fitted without touching either of the other
two. Each shift contributes one decision state per review interval, so a
block of \(n\) shifts yields \(32n\) states: the test corpus of 50
shifts gives \(50 \times 32 = 1600\) states before filtering. On the
labelled test corpus 1525 of 1600 test states survive
(\(\theta = 2.2/|\mathcal{H}| = 0.367\)); 33.4\% of training epochs have
a best/second-best gap below one pooled standard error at \(M=20\). The
filter is applied to the test corpus only, and never to the training
corpus, so no training state is discarded for being difficult.

The horizon is fixed at \(\tau = 4\) for the deployed model; Section 6.4
studies the choice.

By default DAHS uses the soft label above: a state where EEDD and ATC
are near-tied produces a near-uniform target over those two rules, and
the ranker is trained to reproduce that uncertainty rather than to
guess. The arg-max of the cost vector --- a \emph{hard} label --- is the
natural alternative, and Section 6.8 reports an ablation that finds the
two equivalent. We describe the deployed (soft) model here and treat the
label's form as a design choice rather than as a contribution.

The deployed model truncates the rollout at \(\tau = 4\) of up to 32
intervals. We do not claim a bound on the truncation remainder: arrivals
are Poisson, so a finite per-interval cost scale is not proved, and the
implemented window cost is a terminal potential that can be negative.
Horizon is treated as an empirical knob (Section 6.4). Model error is
treated by labelling under nominal dynamics and evaluating frozen
controllers under perturbation (Section 6.11).

\subsection{Regime discovery}\label{regime-discovery}

Warehouse shifts pass through qualitatively distinct operating regimes
--- a quiet opening, a saturated mid-shift, a perishable burst. DAHS
makes this explicit. A Gaussian mixture model is fit to the
training-state features; the number of components is chosen by BIC over
\(K \in \{2,3,4,5,6,7,8,10,12\}\) with five EM restarts per \(K\). On
the regenerated corpus BIC selects \(K^\star = 12\), the \textbf{upper
endpoint of the grid} (mean pairwise ARI \(0.970\), above the \(0.85\)
stability threshold). Because the selected \(K\) sits on the boundary,
the grid --- not a turning point in the data --- chose the number of
regimes; we report that rather than treating twelve as a discovered
structure. The twelve soft regime-membership posteriors are appended to
\(\phi(S_t) \in \mathbb{R}^{26}\), so the ranker sees a
\(38\)-dimensional vector. Regime discovery remains a lightweight
component; Section 6.8 ablates it.

\subsection{The calibrated ranker}\label{the-calibrated-ranker}

The ranker is a gradient-boosted decision-tree classifier (Chen and
Guestrin 2016) with an \(|\mathcal{H}|\)-class soft-probability output
(\(|\mathcal{H}|=6\) after screening). The soft target is fitted by an
inverse-entropy-weighted replication of each training state across the
six classes, which makes the training objective the Kullback--Leibler
divergence between the predicted distribution and the soft label and
down-weights states whose labels are near-uniform (and therefore carry
little discriminative signal). The hard-label variant of Section 6.8
instead uses a standard one-hot cross-entropy; the rest of the pipeline
--- feature set, cross-validation, calibration --- is identical.
Hyperparameters are selected by 5-fold cross-validation with folds
grouped by shift, so no shift contributes states to both a training and
a validation fold. The reference point for the cross-validated soft
cross-entropy is the uniform label, \(\log|\mathcal{H}|\), which moves
with the screened pool. The selected configuration is
\texttt{max\_depth} = 6, \texttt{n\_estimators} = 200,
\texttt{learning\_rate} = 0.03, with mean grouped-CV soft cross-entropy
\(0.875\) against the uniform baseline \(\log 6 \approx 1.79\).

Tree ensembles are not probability-calibrated out of the box. DAHS
post-processes the ranker output with isotonic regression fit on a
held-out 20\% of training shifts. The design target on expected
calibration error is 0.05; Section 6.6 reports the achieved value.
Missing the bar does not change the deployed model.

\subsection{The switching controller}\label{the-switching-controller}

At deployment the calibrated ranker emits, each interval, a distribution
over the retained pool \(\mathcal{H}\). A thin \emph{switching
controller} maps that distribution to an action. It (i) applies the same
expiry-rule mask used at labelling time, which is a no-op on the
screened pool (Section 4.3); (ii) enforces a minimum dwell of
\(T_{\min} = 2\) intervals, counting the switch epoch
(\texttt{\_dwell\_remaining\ =\ t\_min\ -\ 1}), to prevent operationally
disruptive rule thrashing; and (iii) overrides the dwell and switches
immediately when ranker entropy is below half the maximum
(\(H < 0.5\log|\mathcal{H}|\), \(\approx 0.90\) nats at \(K=6\)). We
deliberately frame the controller as a \emph{stability and
constraint-enforcement guardrail}, not a performance driver --- and
Section 6.8 reports, honestly, that removing it slightly \emph{improves}
the headline KPIs. Its role is to make the policy deployable (bounded
switching), not to win the comparison; the ablation quantifies the small
KPI price of that guardrail.

\begin{center}\rule{0.5\linewidth}{0.5pt}\end{center}

\section{Experimental Setup}\label{experimental-setup}

\textbf{Test shifts}. All methods are evaluated on the same 50 held-out
shift seeds, disjoint from the 250 training shifts and fixed once. Every
reported KPI is a mean over these 50 shifts.

\textbf{Scenarios}. Beyond the default operating point, three scenarios
stress the method: \emph{low load} (arrival rate \(1.0\), no
perishables), \emph{balanced} (rate \(1.5\)), and
\emph{high-load-perishable} (rate \(2.2\), perishable probability
\(0.4\), customer windows scaled to \((12,36,72)\) minutes ---
\(0.8\times\) the default triangular). Scenario parameters were fixed
before evaluation and are not tuned per method.

\textbf{Baselines}. We compare DAHS against the static rules retained by
the screen of Section 3.6; \textbf{snapshot\_xgb}, an ablation identical
to DAHS but with the rollout horizon collapsed to \(\tau = 1\),
isolating the value of the horizon; \textbf{LinUCB} (Li et al. 2010), a
contextual bandit, with features standardised (LinUCB updates on the
test shifts and is therefore not a frozen peer); \textbf{PPO} (Schulman
et al. 2017) at 8000 training timesteps (\texttt{total\_timesteps}; not
review epochs), with the hyperparameter sensitivity analysis of Section
6.9 (not a matched simulation budget: labelling uses millions of
interval-steps); and \textbf{offline\_fqi}, a faithful offline
reinforcement-learning competitor --- fitted Q-iteration (Ernst et al.
2005) with FEFO action masking --- a documented no-op when FEFO is
absent from the pool --- an instance of the maskable-action-value family
of Offline-LD (Remmerden et al. 2025). It trains on the same logged
shifts as DAHS, under the same behaviour policy and per-interval reward,
and uses the same gradient-boosted-tree model class and the same feature
set as the DAHS ranker, so that the comparison isolates the training
signal --- a directly measured per-rule cost vector against a single
bootstrapped value --- from the function approximator and from the state
representation. Section 6.10 analyses it, together with the
action-coverage diagnostics that determine whether the comparison is
clean.

\textbf{The teacher: rolling-horizon rollout MPC}. DAHS is trained to
reproduce a \(\tau\)-step rollout, so the controller that simply
\emph{runs} that rollout online is the natural reference.
\textbf{rolling\_mpc} evaluates every rule over \(\tau\) intervals,
averaged over independent continuations that use the same scoring
function as Section 4.3 (not the labelling RNG; teacher evaluation uses
a separate seed offset), commits the arg-min rule for one interval,
discards the remainder of the plan, and replans. Labelling uses \(M=20\)
continuations; the online teachers use \(M=5\) so that a 50-shift
evaluation is affordable. The scoring rule is the same potential
difference, but the Monte Carlo budget is not, so a gap between
rolling\_mpc and DAHS mixes function-approximation error with estimator
noise and cannot be read as a pure distillation gap. Raising the online
teachers to the labelling budget \(M=20\) does not close it:
rolling\_mpc moves from \(363.42\) to \(363.76\) and greedy\_mpc is
unchanged at \(356.98\). At non-terminal epochs a \(\tau=1\) step never
admits \((t,t+L]\), so greedy\_mpc is invariant to \(M\). At the last
epoch, with \texttt{terminal\_admit} on, \texttt{run\_with\_policy} then
admits the resampled tail; the reported \(M\)-invariance of greedy\_mpc
is empirical. greedy\_mpc beating rolling\_mpc on default \(J\) is
consistent with a noisy \(M=5\) arg-min at \(\tau=4\). The one-step
controller is retained as its \(\tau = 1\) special case.

This baseline is what makes the paper's central claim falsifiable.
\emph{How much does distillation lose?} The gap between rolling\_mpc and
DAHS at the same \(\tau\) is the price of replacing a lookahead with a
forward pass. \emph{What does it buy?} We report per-decision latency
for both, so the amortisation appears as a measured ratio rather than an
argument. \emph{Can the student beat the teacher?} Section 6.2 reports
that it does not: both online lookaheads are cheaper, with paired
intervals excluding zero. A fitted selector could in principle
regularise noisy rollouts; that is not what happens here. \emph{And is
the horizon the mechanism?} Comparing rolling\_mpc at \(\tau = 1\)
against \(\tau = 4\) separates the value of lookahead depth from the
value of learning.

\textbf{Objective and metrics}. Every learned method optimises the
composite cost \(J\) of Section 3.3, and \(J\) is the primary metric of
comparison. The breach count is only one of the four terms in the
objective. Making \(J\) primary removes a reporting incentive: a method
could improve one component at the expense of another that went
unreported.

That every learned method is scored against the same \(J\) is the
evaluation contract, not a claim that every arm sees the same
last-interval increment. The objective is defined once, in a single
module. Rollout labels are truncated-horizon estimates of that \(J\) ---
truncation is the method, not a defect to be closed.
\texttt{sim.terminal\_admit} gates DAHS \texttt{run\_with\_policy} and
the static/teacher eval loop. PPO and fitted Q-iteration always fold
leftover arrivals into the last-step potential; they do not read the
flag. LinUCB's payoff is \texttt{env.potential()} after the \emph{next}
epoch's \texttt{observe()}, so it includes newly admitted arrivals, and
\texttt{reset()} between shifts drops the last interval. Frozen DAHS
labels were generated with the flag off and omit last-window leftovers;
live \(|A|\approx 791\) tables align with FQI/PPO on that interval, not
with those labels. The static rules optimise nothing and are evaluated
against \(J\). The objective is written once, in a single module shared
by the labeller and the evaluation harness.

Alongside \(J\) we report its decomposition, so no component can hide:
the \textbf{service-failure rate} --- the share of \emph{arrived} orders
that ship late or spoil, whether or not they were ever dispatched ---
together with mean tardiness, spoilage rate, throughput, unserved and
rejected demand, and picker utilisation. We also report the breach rate
over completed orders only, under that explicit name. Uncertainty is
quantified by 10,000-resample bootstrap 95\% confidence intervals.
\textbf{Confirmatory test for paired composite cost.} A warehouse pays
total cost, not median cost, so the confirmatory functional is the mean.
A method is declared different from DAHS on \(J\) when the paired 95\%
bootstrap interval of the mean difference excludes zero. Wilcoxon
signed-rank \(p\)-values with Benjamini--Hochberg control are a
sensitivity analysis on the sign and median of the paired differences
and are not the confirmatory rule. They are reported the same way for
every claim, not chosen per claim. Several paired differences on this
corpus are tail-driven (Always-EEDD: 21 wins, 7 losses, 22 exact ties;
\(\tau=2\): median difference 0). A mean interval can therefore exclude
zero when the median does not, and the reverse. When that happens we
still follow the mean interval, and we label the claim as a statement
about mean cost. Service-failure rate uses the same interval rule when
it is the claim being tested.

\section{Results}\label{results}

\subsection{Rule calibration, screening, and
complementarity}\label{rule-calibration-screening-and-complementarity}

If one rule were best everywhere, selection would be pointless.
Establishing that it is not requires calibrated rules, a screen that
distinguishes redundant rules from specialists, and evidence of
complementarity across the \textbf{state space} rather than across
instances.

This section is complete: it runs on the 30-shift calibration block,
which is disjoint from both training and test, and it does not depend on
the trained selector. Its results therefore stand ahead of the rest of
Section 6.

\subsubsection{Calibration}\label{calibration}

\textbf{Table 2}. Look-ahead scale calibration, 30 calibration shifts,
\(M = 5\) continuations under common random numbers. Standalone cost is
the composite objective with the rule used alone; portfolio marginal is
the cost increase when the rule is removed from the pool at that \(k\).

{\def\LTcaptype{none} % do not increment counter
\begin{longtable}[]{@{}
  >{\raggedright\arraybackslash}p{(\linewidth - 10\tabcolsep) * \real{0.1304}}
  >{\raggedleft\arraybackslash}p{(\linewidth - 10\tabcolsep) * \real{0.1739}}
  >{\raggedleft\arraybackslash}p{(\linewidth - 10\tabcolsep) * \real{0.1739}}
  >{\raggedleft\arraybackslash}p{(\linewidth - 10\tabcolsep) * \real{0.1739}}
  >{\raggedleft\arraybackslash}p{(\linewidth - 10\tabcolsep) * \real{0.1739}}
  >{\raggedleft\arraybackslash}p{(\linewidth - 10\tabcolsep) * \real{0.1739}}@{}}
\toprule\noalign{}
\begin{minipage}[b]{\linewidth}\raggedright
Rule
\end{minipage} & \begin{minipage}[b]{\linewidth}\raggedleft
\(k^\star_\text{standalone}\)
\end{minipage} & \begin{minipage}[b]{\linewidth}\raggedleft
cost at \(k^\star\)
\end{minipage} & \begin{minipage}[b]{\linewidth}\raggedleft
\(k^\star_\text{portfolio}\)
\end{minipage} & \begin{minipage}[b]{\linewidth}\raggedleft
marginal at \(k^\star\)
\end{minipage} & \begin{minipage}[b]{\linewidth}\raggedleft
deployed
\end{minipage} \\
\midrule\noalign{}
\endhead
\bottomrule\noalign{}
\endlastfoot
ATC & 1.5 & 459.2 & 3.0 & 0.92 & \textbf{3.0} \\
COVERT & 4.0 & 404.3 & 4.0 & 3.21 & \textbf{4.0} \\
\end{longtable}
}

The portfolio value is deployed, because the pool is what ships; both
are reported because ATC is also a standalone benchmark. The Always-ATC
row of the default comparison (Section 6.2) uses the portfolio scale
\(k=3.0\), not the standalone \(k^\star=1.5\). A standalone Always-ATC
at \(k=1.5\) has now been evaluated on the same 50 test shifts and costs
\(J=526.99\), against \(560.76\) for the portfolio scale \(k=3.0\):
refitting \(k\) recovers \(33.8\) of composite cost, but Always-ATC at
its own optimum remains far above DAHS (\(382.27\)).

\textbf{This settles the WSPT/ATC inversion.} ATC's standalone cost is
U-shaped in \(k\) with a minimum of 459.2 at \(k^\star = 1.5\), and
rises monotonically thereafter to 1004.2 at \(k = 20\) --- a factor of
2.19. Since WSPT is exactly the \(k \to \infty\) limit (Section 3.6),
that curve \emph{is} the ATC-to-WSPT interpolation, and it says a fitted
ATC beats WSPT by more than two-fold on this problem. An unfitted
\(k=2.0\) inverts that ordering. The two fitted values also differ by a
factor of two (\(k^\star_\text{standalone} = 1.5\) against
\(k^\star_\text{portfolio} = 3.0\)), which is the concrete case for
calibrating both: the scale that makes ATC best \emph{alone} is not the
scale that makes it most useful \emph{inside a pool}, where its job is
to cover states the other rules handle badly.

\subsubsection{Screening}\label{screening}

\textbf{Table 3}. Pool screening on the same corpus. Screening is by
marginal contribution to composite cost. Marginal contribution is the
increase in achievable cost when the rule is removed, with a percentile
bootstrap interval; a rule is retained when that interval excludes zero.

{\def\LTcaptype{none} % do not increment counter
\begin{longtable}[]{@{}lrrlc@{}}
\toprule\noalign{}
Rule & win rate & marginal contribution & 95\% CI & retained \\
\midrule\noalign{}
\endhead
\bottomrule\noalign{}
\endlastfoot
\textbf{EEDD} & 0.650 & 5.403 & {[}4.840, 5.991{]} & yes \\
COVERT & 0.145 & 2.047 & {[}1.613, 2.527{]} & yes \\
MS & 0.070 & 0.248 & {[}0.119, 0.412{]} & yes \\
ATC & 0.055 & 0.086 & {[}0.032, 0.155{]} & yes \\
MDD & 0.011 & 0.039 & {[}0.014, 0.070{]} & yes \\
EDD & 0.068 & 0.007 & \([8.7{\times}10^{-7}, 0.021]\) & yes \\
FIFO & 0.001 & 0.000 & {[}0.000, 0.000{]} & no \\
WSPT & 0.000 & 0.000 & {[}0.000, 0.000{]} & no \\
FEFO & 0.000 & 0.000 & {[}0.000, 0.000{]} & no \\
\end{longtable}
}

\textbf{FIFO earns nothing.} It is the cost-minimising rule at 0.1\% of
decisions and its marginal contribution is identically zero. As the
zero-information control in a due-date-driven setting that is the
expected result, and it is the direct answer to the question of what
FIFO was doing in the pool: nothing, and it is dropped. Cause 2 in
Section 6.2 is why an arrival-sorted rule can look cheap under
non-causal admission: every arrival-agnostic rule is penalised and FIFO
is not.

\textbf{WSPT earns nothing either, and that is consistent rather than
surprising.} A calibrated ATC dominates it by construction, so once ATC
is fitted, WSPT is a strictly worse member of the same family and
contributes nothing at the margin.

\textbf{FEFO's failure is not evidence against expiry-awareness.} FEFO
ranks on \(x_o\), which is infinite for the 80\% of orders that are not
perishable, so it sorts every non-perishable order behind every
perishable one. It contributes nothing because it is a bad rule on this
order mix, not because the product clock is uninformative --- and
reading its zero as ``expiry-awareness buys nothing'' would have been
the wrong inference. EEDD, which reads \(\min(d_o, x_o)\), settles it:
it wins 65\% of decisions with a marginal contribution of 5.403, two and
a half times COVERT's and by far the largest in the pool.
Expiry-awareness matters a great deal on this problem; the screen was
rejecting a poor implementation of it.

\textbf{EEDD wins 65\% of decisions.} That is above a 0.60 concentration
mark we wrote down before looking at the screen. We did not drop the
rule to get under the mark, and we do not treat the mark as a gate that
changed the rest of the campaign: the question this paper asks is what
distillation is worth on the screened pool, and EEDD is the pool's
dominant member. A pool this concentrated leaves a selector little room.
Whether that room is enough on \(J\) is what Section 6.2 measures, and
the per-cell oracle gap below bounds it.

\subsubsection{Complementarity}\label{complementarity}

Win rate over a grid of the two state dimensions that govern the
decision (queue length and deadline pressure, mean slack), in quantile
bins, is the test of complementarity. A pool is complementary when
different rules own different cells. Figure 1 reports that grid.

\textbf{Table 4}. Cell ownership and the oracle gap over the 4x4 grid
(960 decision epochs from the calibration corpus).

{\def\LTcaptype{none} % do not increment counter
\begin{longtable}[]{@{}lr@{}}
\toprule\noalign{}
Quantity & Value \\
\midrule\noalign{}
\endhead
\bottomrule\noalign{}
\endlastfoot
Grid cells & 16 \\
Cells owned by EEDD & 15 \\
Cells owned by COVERT & 1 \\
Decisions won by the best single rule (always EEDD) & 65.00\% \\
Decisions won by the per-cell oracle & 72.29\% \\
\textbf{Oracle gap over the best single rule} & \textbf{7.29 pp} \\
\end{longtable}
}

\textbf{We report this against ourselves.} One rule owns fifteen of
sixteen cells. At this resolution the pool is not complementary, and a
selector that reads only queue length and deadline pressure could beat
``always EEDD'' on at most 7.29 percentage points of win rate. That is a
small opening, and it is the correct place to say so, before Section 6.2
rather than after it.

Two qualifications matter for reading it, and neither rescues the number
so much as bound it in the other direction.

\emph{The grid oracle is a floor on the available headroom, not a
ceiling on DAHS.} 72.29\% is what a selector restricted to a coarse 4x4
partition of two state dimensions could achieve. DAHS reads 26 features
(Section 4.2), so a finer partition of the state space can only raise
the oracle. The gap is therefore a lower bound on the room available to
a richer selector, and the honest statement is that \emph{this
projection} of the state space does not by itself justify selection.

\emph{Win rate is not the objective.} A rule can win rarely and still be
worth its slot if the states it wins are expensive ones, which is
exactly why Table 3 screens on marginal contribution rather than on win
rate --- and the two orderings differ: COVERT wins 14.5\% of decisions
but carries a marginal contribution of 2.047, while EDD wins 6.8\% and
carries 0.007. The quantity that decides whether selection pays is
composite cost, and Section 6.2 measures it. We did not compute a
composite-cost oracle on the same 4x4 grid. What we can say from the
test set is that the static champion on \(J\) is Always-COVERT, not
Always-EEDD, and that DAHS beats Always-COVERT on 49 of 50 shifts
(Section 6.2). Win-rate concentration on EEDD therefore overstates how
little room there is on the objective.

\begin{figure}
\centering
\pandocbounded{\includegraphics[keepaspectratio,alt={Rule complementarity over the state space: win rate of each retained rule across a grid of queue length (quantile bins) against deadline pressure (mean slack, quantile bins). Complementarity means different rules own different cells.}]{../../figures/S1_calibration/diversity_state_grid.png}}
\caption{Rule complementarity over the state space: win rate of each
retained rule across a grid of queue length (quantile bins) against
deadline pressure (mean slack, quantile bins). Complementarity means
different rules own different cells.}
\end{figure}

Figure 1 is the state-space win-rate grid described above.

\subsection{Main comparison}\label{main-comparison}

Rank the table by composite cost. Table 5 is the live comparison under
the objective of Section 3.3 and causal periodic-review admission.

\subsubsection{Two modelling defects that would invalidate a rule
comparison}\label{two-modelling-defects-that-would-invalidate-a-rule-comparison}

If the objective ignores priority weights, or if the dispatcher idles
pickers on orders that have not arrived, a rule comparison does not
measure rule quality.

\textbf{Cause 1: the objective did not measure what the rules optimise}.
WSPT and ATC rank by \(w_o/p_o\), using the priority weights of Section
3.1. An unweighted objective grades those rules against a criterion they
were not designed for. Section 3.3 puts \(w_o\) into \(J\), so rule and
objective agree.

\textbf{Cause 2: the dispatcher idled pickers on behalf of
arrival-agnostic rules}. Admitting every order that will arrive before
the end of the current interval, and assigning start times
\(\max(\text{picker free}, a_o, t)\), reserves a picker for an order
that is not yet in the building. FIFO, sorted by arrival, never triggers
this; WSPT and ATC, which are arrival-agnostic, do. Section 3.4 admits
only orders that have arrived by \(t\).

All results in this section use that environment and the screened pool
of Section 3.6.

\textbf{Table 5.} Default scenario, 50 test shifts, composite objective
and causal admission. Ranked by composite cost, the quantity every
learned method optimises. Mean arrived \(=791\) (the per-shift count is
not constant; it is identical across methods on each seed). Dropped
\(=0\). Paired 95\% intervals are bootstrap percentile intervals of
(method - DAHS) composite cost over the 50 aligned shifts, 10,000
resamples. The × column is a ratio of means; the bracket is that
difference interval, not a confidence interval for the ratio. Tardiness
is the KPI mean of censored lateness (\texttt{tardiness\_accounted}):
served orders at finish time, unserved orders at \(T\), \textbf{without}
the \(+p_o\) the objective uses in \(J\). The two quantities are not
interchangeable.

{\def\LTcaptype{none} % do not increment counter
\begin{longtable}[]{@{}
  >{\raggedright\arraybackslash}p{(\linewidth - 18\tabcolsep) * \real{0.0789}}
  >{\raggedleft\arraybackslash}p{(\linewidth - 18\tabcolsep) * \real{0.1053}}
  >{\raggedleft\arraybackslash}p{(\linewidth - 18\tabcolsep) * \real{0.1053}}
  >{\raggedleft\arraybackslash}p{(\linewidth - 18\tabcolsep) * \real{0.1053}}
  >{\raggedleft\arraybackslash}p{(\linewidth - 18\tabcolsep) * \real{0.1053}}
  >{\raggedleft\arraybackslash}p{(\linewidth - 18\tabcolsep) * \real{0.1053}}
  >{\raggedleft\arraybackslash}p{(\linewidth - 18\tabcolsep) * \real{0.1053}}
  >{\raggedleft\arraybackslash}p{(\linewidth - 18\tabcolsep) * \real{0.1053}}
  >{\raggedright\arraybackslash}p{(\linewidth - 18\tabcolsep) * \real{0.0789}}
  >{\raggedleft\arraybackslash}p{(\linewidth - 18\tabcolsep) * \real{0.1053}}@{}}
\toprule\noalign{}
\begin{minipage}[b]{\linewidth}\raggedright
Method
\end{minipage} & \begin{minipage}[b]{\linewidth}\raggedleft
Composite cost
\end{minipage} & \begin{minipage}[b]{\linewidth}\raggedleft
SFR
\end{minipage} & \begin{minipage}[b]{\linewidth}\raggedleft
Spoil
\end{minipage} & \begin{minipage}[b]{\linewidth}\raggedleft
Tardiness
\end{minipage} & \begin{minipage}[b]{\linewidth}\raggedleft
Thru.
\end{minipage} & \begin{minipage}[b]{\linewidth}\raggedleft
Util.
\end{minipage} & \begin{minipage}[b]{\linewidth}\raggedleft
Latency (ms)
\end{minipage} & \begin{minipage}[b]{\linewidth}\raggedright
Cost vs DAHS
\end{minipage} & \begin{minipage}[b]{\linewidth}\raggedleft
DAHS wins
\end{minipage} \\
\midrule\noalign{}
\endhead
\bottomrule\noalign{}
\endlastfoot
greedy\_mpc (\(\tau=1\), \(M=5\)) & 356.98 & 0.0590 & 0.0361 & 0.705 &
732.5 & 0.956 & 594 & 0.93× {[}-39.4, -12.8{]} & 19 \\
rolling\_mpc (\(\tau=4\), \(M=5\)) & 363.42 & 0.0624 & 0.0381 & 0.748 &
732.4 & 0.956 & 670 & 0.95× {[}-28.1, -10.8{]} & 15 \\
\textbf{DAHS} & \textbf{382.27} & \textbf{0.0669} & \textbf{0.0384} &
\textbf{0.754} & \textbf{731.9} & \textbf{0.956} & \textbf{4.24} & --- &
--- \\
snapshot\_xgb (\(\tau=1\)) & 388.97 & 0.0651 & 0.0327 & 0.714 & 732.0 &
0.956 & 3.65 & 1.02× {[}-3.6, 18.2{]} & 32 \\
offline\_fqi & 397.64 & 0.0702 & 0.0452 & 0.779 & 732.0 & 0.956 & 3.26 &
1.04× {[}2.9, 29.8{]} & 31 \\
COVERT & 455.21 & 0.0812 & 0.0830 & 0.788 & 732.1 & 0.956 &
\textless0.01 & 1.19× {[}60.9, 85.0{]} & 49 \\
LinUCB & 552.39 & 0.0838 & 0.0672 & 0.751 & 731.2 & 0.956 & --- & 1.45×
{[}94.7, 264.2{]} & 46 \\
ATC (\(k=3.0\)) & 560.76 & 0.1075 & 0.0914 & 1.133 & 734.3 & 0.956 &
\textless0.01 & 1.47× {[}150.9, 206.6{]} & 50 \\
ppo\_fair (untuned) & 611.77 & 0.0924 & 0.0887 & 0.852 & 730.5 & 0.956 &
0.32 & 1.60× {[}114.0, 380.5{]} & 50 \\
EEDD & 696.62 & 0.0915 & 0.0348 & 0.763 & 729.8 & 0.956 & \textless0.01
& 1.82× {[}141.4, 528.8{]} & 21 \\
MDD & 734.00 & 0.0970 & 0.0986 & 0.635 & 730.5 & 0.956 & \textless0.01 &
1.92× {[}200.9, 540.9{]} & 49 \\
EDD & 763.90 & 0.0985 & 0.1000 & 0.725 & 729.8 & 0.956 & \textless0.01 &
2.00× {[}211.9, 594.5{]} & 49 \\
MS & 790.67 & 0.1010 & 0.1024 & 0.760 & 729.1 & 0.956 & \textless0.01 &
2.07× {[}225.1, 635.7{]} & 49 \\
WSPT & 1216.54 & 0.0957 & 0.0591 & 5.387 & 743.3 & 0.955 & \textless0.01
& 3.18× {[}753.4, 911.9{]} & 50 \\
FIFO & 1486.82 & 0.1837 & 0.0720 & 1.963 & 730.2 & 0.956 & \textless0.01
& 3.89× {[}895.5, 1342.4{]} & 50 \\
FEFO & 1699.80 & 0.1971 & 0.0001 & 2.738 & 730.3 & 0.956 & \textless0.01
& 4.45× {[}1077.5, 1584.8{]} & 50 \\
\end{longtable}
}

Wins are shifts on which DAHS is strictly cheaper. Under the
confirmatory paired interval, every static rule, both teachers,
offline\_fqi, LinUCB and Table 5's \texttt{ppo\_fair} row reject
equality with DAHS on composite cost. snapshot\_xgb does not: the
interval includes zero (\([-3.6, 18.2]\)). Wilcoxon+BH on the same 15
comparisons rejects snapshot\_xgb (\(p_{\mathrm{adj}}=0.029\); 32 wins,
14 losses, 4 ties, zeros discarded). On service-failure rate the
snapshot pair is a null (\(p_{\mathrm{adj}}=0.33\)). We treat \(\tau=1\)
as not demonstrably different from deployed DAHS on \(J\) under the
confirmatory rule, and as a null on SFR.

The test block is a weak instrument wherever paired differences are
frequently zero or tail-driven. Always-EEDD is bit-identical to DAHS on
22 of 50 shifts (21/7/22 wins-losses-ties; median paired difference 0).
Under low load (Table 6) thirteen methods share \(J=9.82\). Against
Always-EEDD the paired-difference SD is 705, so an 80\% power minimum
detectable mean effect on this \(n=50\) is about 285 cost units, or 75\%
of DAHS mean \(J\); a 4\% effect has essentially no power. Against
Always-COVERT the same 50 shifts detect a 4.7\% mean effect at 80\%
power (49/50 DAHS wins). Against fitted Q the observed 4\% gap has
approximate paired-\(t\) power \(0.58\); the 80\% minimum detectable
effect is 5.2\% of mean \(J\). Nulls on this block are unsurprising when
many shifts cannot distinguish the methods. They are not a demonstration
that a 4\% effect is absent.

\begin{figure}
\centering
\pandocbounded{\includegraphics[keepaspectratio,alt={Default-scenario composite cost by method, 50 held-out shifts. Paired against DAHS.}]{../../figures/E2/default_forest_composite_cost.png}}
\caption{Default-scenario composite cost by method, 50 held-out shifts.
Paired against DAHS.}
\end{figure}

Figure 2 is that forest plot.

Three facts about the live table:

\begin{enumerate}
\def\labelenumi{\arabic{enumi}.}
\tightlist
\item
  \textbf{The teacher does not beat one-step lookahead}, and DAHS does
  not recover the teacher. Distillation is an amortisation of a
  \emph{worse} scoring rule at \(M=5\), not of an oracle. Per-decision
  latency is the quantity DAHS wins (4.2 ms vs 670 ms for rolling\_mpc;
  Section 6.12).
\item
  \textbf{Always-COVERT, not Always-EEDD, is the static to beat on
  cost.} Win rate on the Section 6.1 grid is the wrong proxy for the
  objective. Against COVERT, DAHS wins 49 of 50 shifts at 1.19 times;
  against EEDD the mean ratio is 1.82 times but DAHS is strictly cheaper
  on only 21 shifts, EEDD on 7, and they tie on 22. The EEDD mean is
  tail-driven (worst shift \(5300\) against DAHS \(2440\)). DAHS also
  wins the median (\(44.2\) against \(57.7\)).
\item
  \textbf{WSPT has the highest throughput} (743 vs FIFO 730) and every
  method sits at picker utilisation \(\approx 0.956\). Cause 2 is the
  mechanism that would have produced idle pickers under non-causal
  admission.
\end{enumerate}

\textbf{Table 6.} Composite cost and service-failure rate across four
scenarios, 50 shifts each. Best static is the cheapest always-on rule in
that scenario.

{\def\LTcaptype{none} % do not increment counter
\begin{longtable}[]{@{}
  >{\raggedright\arraybackslash}p{(\linewidth - 10\tabcolsep) * \real{0.1500}}
  >{\raggedright\arraybackslash}p{(\linewidth - 10\tabcolsep) * \real{0.1500}}
  >{\raggedright\arraybackslash}p{(\linewidth - 10\tabcolsep) * \real{0.1500}}
  >{\raggedright\arraybackslash}p{(\linewidth - 10\tabcolsep) * \real{0.1500}}
  >{\raggedleft\arraybackslash}p{(\linewidth - 10\tabcolsep) * \real{0.2000}}
  >{\raggedleft\arraybackslash}p{(\linewidth - 10\tabcolsep) * \real{0.2000}}@{}}
\toprule\noalign{}
\begin{minipage}[b]{\linewidth}\raggedright
Scenario
\end{minipage} & \begin{minipage}[b]{\linewidth}\raggedright
DAHS \(J\) / SFR
\end{minipage} & \begin{minipage}[b]{\linewidth}\raggedright
Best static
\end{minipage} & \begin{minipage}[b]{\linewidth}\raggedright
Teachers (\(J\))
\end{minipage} & \begin{minipage}[b]{\linewidth}\raggedleft
FQI \(J\)
\end{minipage} & \begin{minipage}[b]{\linewidth}\raggedleft
PPO \(J\)
\end{minipage} \\
\midrule\noalign{}
\endhead
\bottomrule\noalign{}
\endlastfoot
default & 382.27 / 0.0669 & COVERT 455.21 & greedy 356.98; rolling
363.42 & 397.64 & 611.77 \\
high-load-perish & 11430 / 0.542 & \textbf{WSPT 11171} & greedy 11467;
rolling 11341 & 11471 & 23226 \\
balanced & 17.73 / 0.00373 & EEDD 17.62 & greedy 17.92; rolling 17.33 &
18.50 & 35.65 \\
low load & 9.82 / 0.00327 & COVERT (13-way tie at 9.82) & 9.82 & 9.82 &
9.82 \\
\end{longtable}
}

DAHS does not dominate across regimes. Under high-load-perishable WSPT
is cheaper (paired cost difference \(-258\), interval \([-332,-187]\),
BH-reject); Always-ATC is also cheaper than DAHS there (\(J=11{,}380\)).
Under balanced, EEDD is cheaper by \(0.10\) with a null test (interval
\([-0.31, 0.00]\)). Under low load thirteen methods, including DAHS and
both teachers, return identical \(J=9.82\). The default-scenario ranking
is therefore not a licence to write that selection beats any single rule
in every operating region.

\subsubsection{Boundary conditions: what the selector actually does
under
saturation}\label{boundary-conditions-what-the-selector-actually-does-under-saturation}

Calling a lost cell ``a saturation effect'' is an attribution, not an
analysis. There are two candidate explanations and they have different
remedies.

The first is that the selector stops selecting. We measure this as the
\textbf{exponentiated entropy of the deployed-rule distribution}, which
equals \(|\mathcal{H}|\) when the selector spreads across the pool and
1.0 when it has collapsed. On 1,600 default epochs the EEDD share is
0.849 and the COVERT share is 0.138 (switch rate 0.039). Under
high-load-perishable those shares invert: COVERT 0.828, EEDD 0.173.
Under balanced and low load the selector is effectively Always-EEDD
(shares 0.995 and 0.999). Collapse onto a specialist is the mechanism in
the light-load cells; a switch into COVERT is the mechanism under
perishable saturation.

The second is that the guardrail binds. The minimum dwell holds a rule
for \(T_{\min}\) epochs, and under saturation the queue state changes
fastest --- exactly when a stale rule is most costly. We measure this as
the \textbf{blocked-switch rate}: the share of epochs at which the
ranker's arg-max differed from the deployed rule \emph{because the dwell
was still active}. A \(T_{\min}\) sweep \emph{within}
high-load-perishable finds the cost-minimising dwell tied at the
deployed \(T_{\min}=2\) and at \(T_{\min}=3\) (both \(11429.77\));
\(T_{\min}\in\{0,1\}\) is \(11429.85\) and \(T_{\min}=4\) is
\(11432.15\). The guardrail is not the boundary condition in that
scenario.

\subsection{Sample efficiency}\label{sample-efficiency}

Figure 3.

\begin{figure}
\centering
\pandocbounded{\includegraphics[keepaspectratio,alt={Sample efficiency. DAHS mean composite cost versus the number of simulated training shifts, five independent labelling-and-training replicates except n=250 (one replicate, the deployed model). Source figures/data\_efficiency/data\_efficiency\_curve.png.}]{../../figures/data_efficiency/data_efficiency_curve.png}}
\caption{Sample efficiency. DAHS mean composite cost versus the number
of simulated training shifts, five independent labelling-and-training
replicates except \(n=250\) (one replicate, the deployed model). Source
\protect\texttt{figures/data\_efficiency/data\_efficiency\_curve.png}.}
\end{figure}

\textbf{Table 7.} DAHS data-efficiency, composite cost. Budgets
\(\{25,50,100,150,250\}\); five replicates except \(n=250\) (one
replicate by design).

{\def\LTcaptype{none} % do not increment counter
\begin{longtable}[]{@{}
  >{\raggedleft\arraybackslash}p{(\linewidth - 8\tabcolsep) * \real{0.2105}}
  >{\raggedleft\arraybackslash}p{(\linewidth - 8\tabcolsep) * \real{0.2105}}
  >{\raggedleft\arraybackslash}p{(\linewidth - 8\tabcolsep) * \real{0.2105}}
  >{\raggedleft\arraybackslash}p{(\linewidth - 8\tabcolsep) * \real{0.2105}}
  >{\raggedright\arraybackslash}p{(\linewidth - 8\tabcolsep) * \real{0.1579}}@{}}
\toprule\noalign{}
\begin{minipage}[b]{\linewidth}\raggedleft
\(n\) shifts
\end{minipage} & \begin{minipage}[b]{\linewidth}\raggedleft
Replicates
\end{minipage} & \begin{minipage}[b]{\linewidth}\raggedleft
Mean \(J\)
\end{minipage} & \begin{minipage}[b]{\linewidth}\raggedleft
Mean SFR
\end{minipage} & \begin{minipage}[b]{\linewidth}\raggedright
Notes
\end{minipage} \\
\midrule\noalign{}
\endhead
\bottomrule\noalign{}
\endlastfoot
25 & 5 & 449.17 & 0.0716 & One of five collapsed to Always-EEDD
(\(J=696.62\)) \\
50 & 5 & 383.46 & 0.0665 & Already at the deployed level \\
100 & 5 & 383.88 & 0.0669 & \\
150 & 5 & 382.80 & 0.0665 & \\
250 & 1 & 382.27 & 0.0669 & Deployed model \\
\end{longtable}
}

The selector saturates by 50 shifts. The \(n=25\) collapse is a real
failure mode --- not every short corpus yields a working ranker --- but
it is one replicate in five, not the typical outcome. Fitted Q-iteration
on the same budgets is slower to saturate on the pre-admit evaluation
(mean \(J\): 580, 483, 456, 412, 397 at \(n=25,50,100,150,250\); Section
6.10). Those FQI fits were not retrained after terminal admit; Table 5's
live FQI row is \(397.64\).

\subsection{Rollout horizon, and the number of
continuations}\label{rollout-horizon-and-the-number-of-continuations}

The truncation sketch predicts that truncation bias shrinks as \(\tau\)
grows. The operational test is a sweep over \(\tau \in \{1,2,3,4\}\),
each arm labelled at \(M=20\) and used to train an otherwise identical
ranker, evaluated on the same 50 test shifts. Cross-validated soft
cross-entropy on the training labels is not the deployment criterion; we
report it alongside \(J\) because it is what the ranker actually fits.

\textbf{Table 8.} Rollout-horizon sweep. Cost intervals are 95\%
bootstrap intervals of the 50-shift mean. Those intervals are wide
because they are marginal means, not the confirmatory paired test of
Section 5. \(\tau=1\) is cost-identical to snapshot\_xgb on non-timing
KPIs (latency columns differ); \(\tau=4\) is cost-identical to Table 5's
DAHS on non-timing KPIs.

{\def\LTcaptype{none} % do not increment counter
\begin{longtable}[]{@{}
  >{\raggedleft\arraybackslash}p{(\linewidth - 10\tabcolsep) * \real{0.1739}}
  >{\raggedleft\arraybackslash}p{(\linewidth - 10\tabcolsep) * \real{0.1739}}
  >{\raggedright\arraybackslash}p{(\linewidth - 10\tabcolsep) * \real{0.1304}}
  >{\raggedleft\arraybackslash}p{(\linewidth - 10\tabcolsep) * \real{0.1739}}
  >{\raggedleft\arraybackslash}p{(\linewidth - 10\tabcolsep) * \real{0.1739}}
  >{\raggedleft\arraybackslash}p{(\linewidth - 10\tabcolsep) * \real{0.1739}}@{}}
\toprule\noalign{}
\begin{minipage}[b]{\linewidth}\raggedleft
\(\tau\)
\end{minipage} & \begin{minipage}[b]{\linewidth}\raggedleft
\(J\)
\end{minipage} & \begin{minipage}[b]{\linewidth}\raggedright
95\% CI of mean
\end{minipage} & \begin{minipage}[b]{\linewidth}\raggedleft
SFR
\end{minipage} & \begin{minipage}[b]{\linewidth}\raggedleft
Median label entropy
\end{minipage} & \begin{minipage}[b]{\linewidth}\raggedleft
Test rows kept
\end{minipage} \\
\midrule\noalign{}
\endhead
\bottomrule\noalign{}
\endlastfoot
1 & 388.97 & {[}236.3, 567.7{]} & 0.0651 & 0.693 & 915 / 1600 \\
2 & \textbf{373.16} & {[}225.7, 546.6{]} & \textbf{0.0633} & 0.649 &
1273 / 1600 \\
3 & 374.23 & {[}226.0, 549.0{]} & 0.0637 & 0.638 & 1444 / 1600 \\
4 & 382.27 & {[}233.3, 556.1{]} & 0.0669 & 0.638 & 1525 / 1600 \\
\end{longtable}
}

Every \emph{marginal} mean interval overlaps every other; those
intervals are the wrong test for a \(\tau\) claim (they would also fail
to separate DAHS from Always-COVERT). The confirmatory paired interval
of \((\tau-4)\) on the same 50 shifts is: \(\tau=1\) \(+6.71\)
\([\,-3.63,\,18.17\,]\) (includes 0); \(\tau=2\) \(-9.11\)
\([\,-17.61,\,-1.35\,]\) (excludes 0); \(\tau=3\) \(-8.04\)
\([\,-16.90,\,0.28\,]\) (includes 0). \(\tau=2\) is cheaper than the
deployed \(\tau=4\) model on mean \(J\). The same 50 shifts split 21
cheaper / 20 worse / 9 ties; the median paired difference is 0; Wilcoxon
signed-rank \(p=0.221\). That is a mean-driven interval, the same tail
pattern disclosed for Always-EEDD (21/7/22). \(\tau=3\) includes 0 on
the interval; its unadjusted Wilcoxon \(p\) is \(0.0495\) (BH \(m=3\),
\(p_{\mathrm{adj}}=0.074\)). On SFR, \(\tau=1\) versus \(\tau=4\)
includes 0 (\(p_{\mathrm{adj}}=0.33\)). The deployed ranker is
\(\tau=4\) because the labels were frozen there. We do not read Table 8
as confirming an interior \(\tau^\star\). Figure 4.

\begin{figure}
\centering
\pandocbounded{\includegraphics[keepaspectratio,alt={Composite cost versus rollout horizon \textbackslash tau. Source results/E4/tau\_summary.parquet.}]{../../figures/E4/tau_composite_cost.png}}
\caption{Composite cost versus rollout horizon \(\tau\). Source
\protect\texttt{results/E4/tau\_summary.parquet}.}
\end{figure}

The companion sweep is \(M \in \{1,5,10,20,40\}\). All five cells
completed. Figure 5 plots composite cost against \(M\).

\textbf{Table 9.} Number of continuations. \(M=20\) is the deployed
model. Entropy is in-band at every \(M\). \(M=1\) has identically zero
rollout SE, so the fraction of labels with separation below one SE is
zero by construction. \(M=1\) post-calibration ECE \(0.067\) misses the
\(0.05\) bar.

{\def\LTcaptype{none} % do not increment counter
\begin{longtable}[]{@{}
  >{\raggedleft\arraybackslash}p{(\linewidth - 12\tabcolsep) * \real{0.1481}}
  >{\raggedleft\arraybackslash}p{(\linewidth - 12\tabcolsep) * \real{0.1481}}
  >{\raggedleft\arraybackslash}p{(\linewidth - 12\tabcolsep) * \real{0.1481}}
  >{\raggedleft\arraybackslash}p{(\linewidth - 12\tabcolsep) * \real{0.1481}}
  >{\raggedleft\arraybackslash}p{(\linewidth - 12\tabcolsep) * \real{0.1481}}
  >{\raggedleft\arraybackslash}p{(\linewidth - 12\tabcolsep) * \real{0.1481}}
  >{\raggedright\arraybackslash}p{(\linewidth - 12\tabcolsep) * \real{0.1111}}@{}}
\toprule\noalign{}
\begin{minipage}[b]{\linewidth}\raggedleft
\(M\)
\end{minipage} & \begin{minipage}[b]{\linewidth}\raggedleft
\(J\)
\end{minipage} & \begin{minipage}[b]{\linewidth}\raggedleft
SFR
\end{minipage} & \begin{minipage}[b]{\linewidth}\raggedleft
Median entropy
\end{minipage} & \begin{minipage}[b]{\linewidth}\raggedleft
frac \textless1 SE
\end{minipage} & \begin{minipage}[b]{\linewidth}\raggedleft
Test kept
\end{minipage} & \begin{minipage}[b]{\linewidth}\raggedright
ECE pre to post
\end{minipage} \\
\midrule\noalign{}
\endhead
\bottomrule\noalign{}
\endlastfoot
1 & 381.08 & 0.0664 & 0.641 & 0 & 1150 / 1600 & 0.077 -\textgreater{}
0.068 \\
5 & 381.28 & 0.0660 & 0.644 & 0.567 & 1384 / 1600 & 0.123
-\textgreater{} 0.034 \\
10 & 383.25 & 0.0670 & 0.649 & 0.455 & 1457 / 1600 & 0.155
-\textgreater{} 0.026 \\
20 & 382.27 & 0.0669 & 0.638 & 0.334 & 1525 / 1600 & 0.170
-\textgreater{} 0.021 \\
40 & 383.63 & 0.0666 & 0.652 & 0.215 & 1526 / 1600 & 0.189
-\textgreater{} 0.018 \\
\end{longtable}
}

Paired against \(M=20\), \(M=1\) differs by \(-1.18\) in \(J\) (95\% CI
\([-7.83, 5.18]\), Wilcoxon \(p=0.91\)); \(M=5\), \(M=10\) and \(M=40\)
likewise include zero (\(M=40\): \(+1.36\), \([-3.05, 5.99]\),
\(p=0.35\)). The five-cell span on mean \(J\) is about \(0.7\%\).
Multi-sample labels improve calibration ECE and keep more test rows
through the ambiguity filter. They do not improve deployed cost. Why
\(M > 1\) remains the argument of Section 4.3: a single continuation is
a hindsight-optimal path, not an estimate of expected cost. That is not
a claim that \(M=20\) is an accuracy lever.

\begin{figure}
\centering
\pandocbounded{\includegraphics[keepaspectratio,alt={Composite cost versus number of continuations M. Source results/E4/n\_samples\_summary.parquet.}]{../../figures/E4/n_samples_composite_cost.png}}
\caption{Composite cost versus number of continuations \(M\). Source
\protect\texttt{results/E4/n\_samples\_summary.parquet}.}
\end{figure}

\subsection{Robustness across untuned
configurations}\label{robustness-across-untuned-configurations}

\begin{figure}
\centering
\pandocbounded{\includegraphics[keepaspectratio,alt={Robustness grid across 12 untuned configurations (4 arrival rates × 3 SLA tightnesses). Heat map of composite cost. The outlined cell is the Table 5 default and matches it to machine precision.}]{../../figures/E8/robustness_grid_heatmap_composite_cost.png}}
\caption{Robustness grid across 12 untuned configurations (4 arrival
rates × 3 SLA tightnesses). Heat map of composite cost. The outlined
cell is the Table 5 default and matches it to machine precision.}
\end{figure}

Figure 6 is that heat map.

Five methods are frozen across the grid: DAHS, greedy\_mpc,
snapshot\_xgb, Always-EEDD and Always-COVERT. Teachers replan with each
cell's true arrival rate and tightness; DAHS and snapshot\_xgb are the
default-trained rankers with no retraining. greedy\_mpc wins 8 of 12
cells and EEDD wins the four light-load default/loose cells. DAHS wins
none, and adding Always-COVERT does not change that: COVERT also wins
none. The \texttt{arr1.65\_default} cell reproduces Table 5 exactly on
every non-timing column. The grid is a stress test of transfer, not a
second evaluation of the default ranking, and it says the one-step
teacher --- not the distilled student --- is the most robust under load
and tightness changes the ranker was not retrained on.

Cell wins are winner-take-all and understate how the losers lose.
Measured as regret against each cell's own winner, DAHS averages
\(17.2\%\) (worst \(55.2\%\)) while Always-EEDD averages \(55.5\%\)
(worst \(125.8\%\)): EEDD's four wins are all light-load cells where
every method is within noise of a near-zero cost, and it degrades by
more than \(95\%\) in each of the four heaviest cells. DAHS never wins a
cell and never collapses in one; that is a different claim from
robustness-best, and it is the one the grid supports.

\subsection{Calibration and
interpretability}\label{calibration-and-interpretability}

Isotonic calibration on a 20\% held-out shift split improves expected
calibration error from 0.1700 to 0.0213 (design target 0.05) and Brier
score from 0.1730 to 0.1240, and \emph{degrades} soft cross-entropy from
0.828 to 2.358. EDD never wins a label on the calibration split and is
passed through uncalibrated. The reliability diagrams are Figure 7.

\begin{figure}
\centering
\pandocbounded{\includegraphics[keepaspectratio,alt={Reliability diagrams before and after isotonic calibration.}]{../../figures/E5/reliability_pre_post.png}}
\caption{Reliability diagrams before and after isotonic calibration.}
\end{figure}

Global SHAP values (Lundberg and Lee 2017), computed as a
\emph{diagnostic} on labelled test states (\texttt{data/test.parquet}),
on the 38-dimensional ranker input (26-feature observation plus 12
regime posterior coordinates) rank \texttt{queue\_length} (0.370),
\texttt{queue\_length\_lag\_1} (0.152), \texttt{mean\_slack\_minutes}
(0.120), \texttt{interval\_index\_in\_shift} (0.053) and
\texttt{queue\_length\_lag\_2} (0.052). The sixth coordinate is a regime
one-hot (\texttt{regime\_post\_9}, 0.040). Congestion and slack
dominate; the product-clock features do not. The remaining 33 of 38
ranker coordinates are not earning their dimensionality on cost at a
detectable margin on this test set (Section 6.8). That is consistent
with Section 3.5: perishability binds, but the customer clock is the
more frequent pivot, and EEDD already reads both clocks inside the
action. Figure 8.

\begin{figure}
\centering
\pandocbounded{\includegraphics[keepaspectratio,alt={Global SHAP feature importance for the ranker.}]{../../figures/E5/shap_summary.png}}
\caption{Global SHAP feature importance for the ranker.}
\end{figure}

\subsection{Real-data grounding}\label{real-data-grounding}

\textbf{Fitting, not validating}. Inter-arrivals are fitted to a public
trace. Processing time and shelf life are design parameters. Appendix C
reports the candidate families, the AIC table, and the selected family.
Inter-arrivals are lognormal (Olist, \(n=98{,}241\)); the exponential
alternative is worse by \(\Delta\)AIC \(=34{,}502\). Customer windows
are the closest shape match (KS \(D=0.039\)). Processing time is not
fitted: the trace's purchase-to-approval delay is not pick time.
Perishable fraction 0.20 is a design choice (Olist food/drink share is
0.0099). The operating arrival rate 1.65 / min is an operating point,
not a fit. Figure 9.

\begin{figure}
\centering
\pandocbounded{\includegraphics[keepaspectratio,alt={Simulator input distributions against the Olist order trace (mean-normalised).}]{../../figures/A/olist_validation.png}}
\caption{Simulator input distributions against the Olist order trace
(mean-normalised).}
\end{figure}

A frozen-ranker replay under empirical Olist arrivals (burstiness CV
\(2.68\), against Poisson CV \(1.00\)) preserves the ranking:
greedy\_mpc 1784, DAHS 1816, snapshot\_xgb 1840, Always-EEDD 3807.
Absolute costs inflate; the teacher still leads and EEDD still trails.
Figure 10.

\begin{figure}
\centering
\pandocbounded{\includegraphics[keepaspectratio,alt={Method KPIs under Poisson against empirical-Olist bursty arrivals (frozen rankers).}]{../../figures/A2/olist_arrivals_compare.png}}
\caption{Method KPIs under Poisson against empirical-Olist bursty
arrivals (frozen rankers).}
\end{figure}

\subsection{Ablations}\label{ablations}

\textbf{Table 10.} Retrain and inference ablations versus DAHS, 50
paired shifts. Composite cost is the column that decides. The
confirmatory functional is the mean (Section 5). \texttt{no\_regime} is
worse than DAHS on mean \(J\) (\(+2.14\), \([0.04, 4.61]\), excludes 0).
That lower bound is \(0.04\) on a cost of \(382\), and Wilcoxon+BH on
the eight-arm family in \texttt{e3\_summary.parquet} (including
\texttt{random\_ambiguity\_filter}, which is per-shift identical to
DAHS) does not reject (\(p_{\mathrm{adj}}=0.34\)). We report a mean-cost
difference of negligible operational size, not a finding that the regime
layer is load-bearing. The other displayed retrain arms include 0.
\texttt{random\_ambiguity\_filter} is omitted from the displayed rows;
the reported \(p_{\mathrm{adj}}\) values are the 8-arm corrections. A
\texttt{top5\_features} retrain that selected coordinates from a
test-set SHAP ranking is omitted from this table; it is a diagnostic in
the paragraph below, not a held-out ablation. Non-rejection is not an
equivalence test. The training wall-clock to convergence is reported for
retrain arms; inference-only arms do not retrain.

{\def\LTcaptype{none} % do not increment counter
\begin{longtable}[]{@{}
  >{\raggedright\arraybackslash}p{(\linewidth - 10\tabcolsep) * \real{0.1364}}
  >{\raggedleft\arraybackslash}p{(\linewidth - 10\tabcolsep) * \real{0.1818}}
  >{\raggedright\arraybackslash}p{(\linewidth - 10\tabcolsep) * \real{0.1364}}
  >{\raggedleft\arraybackslash}p{(\linewidth - 10\tabcolsep) * \real{0.1818}}
  >{\raggedleft\arraybackslash}p{(\linewidth - 10\tabcolsep) * \real{0.1818}}
  >{\raggedleft\arraybackslash}p{(\linewidth - 10\tabcolsep) * \real{0.1818}}@{}}
\toprule\noalign{}
\begin{minipage}[b]{\linewidth}\raggedright
Ablation
\end{minipage} & \begin{minipage}[b]{\linewidth}\raggedleft
\(J\)
\end{minipage} & \begin{minipage}[b]{\linewidth}\raggedright
\(J\) vs DAHS {[}95\% CI{]}
\end{minipage} & \begin{minipage}[b]{\linewidth}\raggedleft
\(p_{\mathrm{adj}}\)
\end{minipage} & \begin{minipage}[b]{\linewidth}\raggedleft
SFR
\end{minipage} & \begin{minipage}[b]{\linewidth}\raggedleft
Train wall (s)
\end{minipage} \\
\midrule\noalign{}
\endhead
\bottomrule\noalign{}
\endlastfoot
hard\_labels & 383.28 & +1.01 {[}-5.52, 7.13{]} & 0.83 & 0.0665 &
1302 \\
no\_calibration & 382.31 & +0.04 {[}-5.76, 4.89{]} & 0.83 & 0.0665 &
--- \\
no\_regime & 384.41 & +2.14 {[}0.04, 4.61{]} & 0.34 & 0.0671 & 1124 \\
no\_switching\_controller & 381.27 & -0.99 {[}-4.24, 1.96{]} & 0.83 &
0.0667 & --- \\
\texttt{single\_sample\_rollout} (\(M=1\)) & 381.08 & -1.18 {[}-7.83,
5.18{]} & 1.00 & 0.0664 & --- \\
DAHS & 382.27 & --- & --- & 0.0669 & --- \\
\end{longtable}
}

The hard-label comparison is a null on both composite cost and
service-failure rate. Soft labels remain the default because they are
the deployed configuration, not because they outperform one-hot arg-max.
Removing the switching controller slightly \emph{improves}
point-estimate cost; the wrapper is a deployability guardrail paid in
KPI, not a performance component. A \texttt{top5\_features} retrain that
selected five coordinates from a test-set SHAP ranking scored
\(J=383.99\) (\(+1.72\), \([-3.44, 6.59]\)). That arm is not in Table
10. It is a diagnostic with selection on the evaluation corpus, not a
confirmation-set protocol.

A single-sample rollout ablation (\(M=1\) relabel and retrain) is the
\(M=1\) cell of Table 9 and the \texttt{single\_sample\_rollout} row of
Table 10 (cost-identical parquets; latency columns differ). It matches
\(M=20\) on \(J\). The \(\tau=1\) snapshot is Table 8.

\subsection{On the PPO baseline}\label{on-the-ppo-baseline}

Untuned PPO at 8000 timesteps (\texttt{ppo\_baseline}, \texttt{n\_steps}
= 64) collapses to a constant EEDD policy: its 50-shift cost is
696.616200, identical to Always-EEDD to machine precision. That is not
the Table 5 \texttt{ppo\_fair} row (\(J=611.77\)), which is a separately
trained stock Stable-Baselines3 policy with a full KPI record.
Sensitivity analysis: a 12-configuration grid was scored on the 50 test
shifts (\texttt{experiments/rl\_sensitivity.py}). Eight of the twelve
cells later score the same calibration cost \(587.36\), so the
calibration split barely discriminates and cannot repair a test-first
search. We do not treat ``the same cell wins both rankings'' as evidence
that selection was on calibration. Table 5 keeps the untuned
\texttt{ppo\_fair} row, which was declared before this grid. The grid is
a diagnostic of how much normalisation buys, not a replacement baseline.

\textbf{Table 11.} PPO sensitivity grid. Against the Always-EEDD
\emph{collapse} row (\(J=696.62\)), \texttt{gap\_closed\_fraction} \(=\)
\((\text{collapse}-\text{best})/(\text{collapse}-\text{DAHS}) = 0.783\).
Against Table 5's named untuned \texttt{ppo\_fair} row (\(J=611.77\))
the same best cell closes \(0.70\). We do not call both denominators
``untuned.'' Factor spreads: normalisation \(313.5\); \texttt{n\_steps}
\(22.5\); \(\gamma\), GAE \(\lambda\) and entropy coefficient \(0\).

{\def\LTcaptype{none} % do not increment counter
\begin{longtable}[]{@{}lr@{}}
\toprule\noalign{}
Configuration & \(J\) \\
\midrule\noalign{}
\endhead
\bottomrule\noalign{}
\endlastfoot
8 of 12 (including the untuned default) & 696.62 \\
\texttt{n\_steps}=256 & 674.14 \\
obs-norm only & 541.51 \\
\textbf{obs-norm and reward-norm} & \textbf{450.45} \\
reward-norm only & 763.90 (Always-EDD) \\
\end{longtable}
}

Tuning closes a substantial share of the gap. The reading that PPO's
deficit is structural rather than budgetary is withdrawn. The winner
(\texttt{norm(obs=True,\ rew=True)}, \(J=450.45\)) is a sensitivity cell
scored on the test block, not a baseline. Eight of twelve calibration
costs tie at \(587.36\), so a calibration-first pass would barely
discriminate. The \(0.783\) figure is the collapse-to-DAHS fraction;
\(0.70\) is the \texttt{ppo\_fair}-to-DAHS fraction. Table 5 keeps the
untuned \texttt{ppo\_fair} row, which is the policy with a full KPI
record; the grid cell is reported here as a diagnostic.

\subsection{On the offline reinforcement-learning
baseline}\label{on-the-offline-reinforcement-learning-baseline}

\textbf{Action coverage.} The behaviour policy is \texttt{random}.
Round-robin over the pool would not be: the interval index is an
observed feature, so \(a = t \bmod |\mathcal{H}|\) is a deterministic
function of the observation and conditional coverage would be one action
per state by construction. On 8,000 training states the effective number
of actions is 5.999 of 6 overall and 5.995 of 6 on the breach-prone
quartile; the interval-conditional mean is 5.937. Coverage is adequate:
the offline-RL deficit cannot be attributed to unsupported actions in
the hard region.

After an observe-once logger (calling \texttt{observe()} twice per
interval zeros \texttt{n\_arrivals\_last\_interval} on 31 of 32 states),
DAHS beats fitted Q on composite cost, \(382.27\) versus \(397.64\)
(ratio \(1.04\), paired interval \([2.90, 29.83]\), confirmatory
interval excludes 0). The margin is about 4\%. It is the only
learned-baseline comparison in Table 5 that favours the selector.
Section 8 records that it is not a demonstration about value learning.
Figure 11.

\begin{figure}
\centering
\pandocbounded{\includegraphics[keepaspectratio,alt={Sample efficiency: DAHS versus fitted Q-iteration at matched shift budgets. The FQI curve is the pre-admit evaluation (\textbar A\textbar=767 at n=250, J=396.80); Table 5's live FQI row is 397.64.}]{../../figures/E9/data_efficiency_offline_fqi.png}}
\caption{Sample efficiency: DAHS versus fitted Q-iteration at matched
shift budgets. The FQI curve is the pre-admit evaluation (\(|A|=767\) at
\(n=250\), \(J=396.80\)); Table 5's live FQI row is \(397.64\).}
\end{figure}

Frozen across the four scenarios and the twelve-cell grid, FQI tracks
DAHS rather than collapsing. Under high-load-perishable both sit near
\(11{,}430\)--\(11{,}471\) while PPO reaches \(23{,}226\) and
Always-EEDD \(24{,}065\). Under default, FQI is the fifth method in
Table 5, between snapshot\_xgb and Always-COVERT.

\subsection{Model misspecification: labelling in one world, deploying in
another}\label{model-misspecification-labelling-in-one-world-deploying-in-another}

The experiment labels once under the nominal simulator and evaluates
frozen controllers --- DAHS, rolling\_mpc, offline\_fqi, Always-EEDD,
Always-COVERT --- under perturbed dynamics: arrival rate
\(\times\{0.8,0.9,1.0,1.1,1.25\}\), processing time on the same grid,
SLA \(\times\{0.8,1.0,1.25\}\), shelf life \(\times\{0.8,1.0,1.25\}\),
and picker count \(\Delta \in \{-2,-1,0,1,2\}\). The online teacher is
pinned to the \emph{nominal} model, so it replans with a wrong model
rather than with the truth.

\textbf{Table 12.} Mean relative-degradation slope of composite cost
against the perturbation, averaged over axes.

{\def\LTcaptype{none} % do not increment counter
\begin{longtable}[]{@{}lrl@{}}
\toprule\noalign{}
Method & Mean slope & Rank (slower degradation first) \\
\midrule\noalign{}
\endhead
\bottomrule\noalign{}
\endlastfoot
COVERT & 18.7 & most robust (model-free) \\
offline\_fqi & 21.9 & \\
DAHS & 22.7 & \\
rolling\_mpc & 23.6 & \\
EEDD & 27.7 & least robust \\
\end{longtable}
}

DAHS is not the most robust model-based method; FQI's mean slope is
slightly shallower. Always-COVERT, which carries no model, degrades
slowest. Always-EEDD degrades fastest on load, capacity and processing
time. On the SLA axis DAHS has the \emph{steepest} slope of the five.
Shelf-life slopes are near zero.

This experiment does not retrain at \(\tau \in \{1,2,3,4\}\) as
\(\varepsilon\) grows, so it does not test whether an operationally best
horizon would shorten. We report the slopes and that the
\(\tau\)--\(\varepsilon\) interaction is unmeasured.

\subsection{Computational cost, and scaling in the size of the rule
pool}\label{computational-cost-and-scaling-in-the-size-of-the-rule-pool}

\subsubsection{Offline cost}\label{offline-cost}

Walking each shift forward once and branching at each epoch (Section
4.3) costs

\[ N + N \cdot |\mathcal{H}| \cdot M \cdot \tau \quad \text{interval-steps per shift,} \]

linear in \(N\): the labeller walks each shift once and branches at each
epoch, rather than replaying each shift from \(t=0\) for every
candidate. For the deployed setup --- \(N=32\), \(|\mathcal{H}|=6\),
\(\tau=4\), \(M=20\), 250 training shifts --- labelling consumed
4,401,600 interval-steps (3,668,000 train, 733,600 test) in \(758\) s
train and \(164\) s test wall-clock on the campaign machine (Intel Core
i9-14900K, 24 cores / 32 threads, 64 GB, Windows 11, Python 3.12.10).
Rule calibration, which sweeps \(k\) under \(M=5\), consumed 2,966,400
additional interval-steps. A separate single-thread
\texttt{compute\_budget\ measure} pass on three shifts recorded
\(389.3\) interval-steps per second (\(2.75\) h single-core for the
3,848,000-step corpus formula used by that driver). Labelling wall-clock
above is the labeller's own multi-core measurement and remains the
campaign figure.

\subsubsection{Online cost, and per-decision inference
latency}\label{online-cost-and-per-decision-inference-latency}

\textbf{Table 13.} Mean per-decision inference latency, 50 default
shifts. The p95 column is the mean, across shifts, of each shift's
per-decision p95 (so DAHS can have mean 4.24 ms and p95 4.25 ms). It is
not a pooled p95 over all decisions.

{\def\LTcaptype{none} % do not increment counter
\begin{longtable}[]{@{}lrrr@{}}
\toprule\noalign{}
Method & Mean (ms) & p95 (ms) & Wall-clock s / shift \\
\midrule\noalign{}
\endhead
\bottomrule\noalign{}
\endlastfoot
DAHS & 4.24 & 4.25 & 0.163 \\
snapshot\_xgb & 3.65 & 3.52 & 0.143 \\
offline\_fqi & 3.26 & 5.12 & 0.130 \\
ppo\_fair & 0.32 & 0.46 & 0.024 \\
greedy\_mpc & 594 & 859 & 19.0 \\
rolling\_mpc & 670 & 941 & 21.5 \\
COVERT / EEDD / FIFO & \textless1e-3 & \textless1e-3 & 0.01 \\
\end{longtable}
}

DAHS is \(158\) times faster per decision than rolling\_mpc (4.24 ms vs
670 ms) and \(140\) times faster than greedy\_mpc. Static rules remain
three orders of magnitude faster still. The amortisation claim is this
ratio, on named hardware, not an argument that a forward pass is faster
than a tree of simulations. At a 15-minute review interval, 670 ms is
\(0.07\%\) of the epoch, so the latency product is real and
operationally unused unless a tighter review is required.

Labelling the \(M=20\), \(\tau=4\) corpus took \(922\) s (\(15.4\) min)
of wall-clock. Ranker retraining on the same machine takes 15--27
minutes (Table 10: 897--1302 s for the retrain arms). At
\(21.5-0.16=21.3\) extra seconds per shift of online lookahead,
labelling alone is repaid after about \(43\) shifts; labelling plus
retraining is repaid after about \(90\)--\(120\) shifts. Both figures
are hardware-specific and ignore that the teacher is also the better
policy.

\subsubsection{\texorpdfstring{Scaling in
\(|\mathcal{H}|\)}{Scaling in \textbar\textbackslash mathcal\{H\}\textbar{}}}\label{scaling-in-mathcalh}

Labelling cost is linear in \(|\mathcal{H}|\) under uniform allocation
of \(M\) continuations to every rule. Two mitigations are discussed.

Adaptive sample allocation (successive halving) discards clearly
dominated rules part-way through the \(M\) continuations and spends the
remaining budget on the survivors. Hierarchical selection first screens
a cheap one-step score and only then rolls the shortlist to depth
\(\tau\). Successive halving is implemented in the labeller as
\texttt{costs\_at\_epoch\_successive\_halving} and was measured as a
compute-budget diagnostic; production labelling
(\texttt{label\_one\_shift}) never calls it and stays on uniform
\texttt{costs\_at\_epoch}. A 50-shift diagnostic at the deployed
\((M,\tau)\) on the \textbf{nine-candidate screen set} (not the six-rule
deployed pool) produced arg-max agreement \(0.856\) against uniform
allocation, mean label KL \(0.185\), and a \(1.1\%\) step saving
(\texttt{successive\_halving.json}). That KL is between untempered
softmaxes of the two cost vectors, not the production tempered labels
\(T(s)=\beta\hat\sigma(s)\). The verdict is unsuccessful, so production
labels stay uniform. Hierarchical selection is described here and is not
implemented. We did not retrain DAHS at pool sizes \(2\), \(4\) and
\(8\). The statement we can support is the complexity: uniform labelling
is \(\Theta(N|\mathcal{H}|M\tau)\) interval-steps per shift; successive
halving replaces \(|\mathcal{H}|M\) with \(|\mathcal{H}|\log M\) in the
usual approximation when gaps are large relative to rollout SE. At
\(M=20\), \(33\%\) of training epochs have a best-versus-second gap
below one pooled SE (Section 4.3), which is exactly the regime in which
early discarding is least safe.

\begin{center}\rule{0.5\linewidth}{0.5pt}\end{center}

\section{Discussion}\label{discussion}

The campaign supports a smaller claim than the one we started with, and
that is the claim we are willing to defend.

A practitioner who already has a resettable simulator and who can afford
\(670\) ms per decision should run the one-step teacher. A practitioner
who needs a millisecond decision can distil it, and should not expect
the student to catch the teacher or to win the transfer grid. There is
no third branch in which DAHS is the policy to run. If the simulator
does not exist, these labels cannot be generated.

Two comparisons that looked structural are not. PPO's deficit was a
missing normalisation wrapper: the collapse-to-DAHS
\texttt{gap\_closed\_fraction} is \(0.783\); the
\texttt{ppo\_fair}-to-DAHS fraction is \(0.70\). Round-robin FQI
coverage would be degenerate by construction; under a random behaviour
policy coverage is adequate and a 4\% cost gap remains. Both baselines
belong in the table. Neither is a demonstration that value-based or
policy-gradient methods cannot do this task. PPO is not a matched
simulation-step budget. The paper does not rest on the 4\% FQI gap.

Perishability binds at this horizon (Table 1) at the fraction we ran.
FEFO drives spoilage essentially to zero at catastrophic cost (Table 5).
EEDD is the rule that reads both clocks. That does not imply that a
selector over those rules will win every regime: WSPT wins
high-load-perishable, and Always-EEDD is statistically tied with DAHS
under balanced load.

\begin{center}\rule{0.5\linewidth}{0.5pt}\end{center}

\section{Limitations}\label{limitations}

\subsection{Shared-simulator
circularity}\label{shared-simulator-circularity}

Every method in Table 5, teachers included, is labelled or trained in
the same simulator it is evaluated in. The comparison of training
signals is internally valid. It is not a claim about transfer to a
physical warehouse. Section 6.11 perturbs the evaluation dynamics while
keeping the labels nominal; that is a start, not a substitute for a
plant. A practitioner without a trustworthy simulator should prefer
value learning from logs (Section 2.4).

\subsection{A small heuristic pool}\label{a-small-heuristic-pool}

Nine candidates were screened and six retained. Adaptive sample
allocation (successive halving) is implemented in the labeller;
hierarchical selection is described in Section 6.12 and is not
implemented. We did not measure accuracy as a function of pool size. The
labelling cost grows with \(|\mathcal{H}|\). A larger, generated pool
--- genetic programming of the low-level rules (Branke et al. 2016;
Nguyen et al. 2017) followed by the same selector --- is a different
paper. EEDD's 65\% win rate already warns that a small hand-designed
pool can concentrate.

\subsection{A single warehouse
setting}\label{a-single-warehouse-setting}

One layout, ten pickers, one-order tours, exogenous routing, Poisson
default arrivals. Appendix C fits inter-arrival shape to Olist;
processing time and shelf life are design parameters. The perishable
fraction 0.20 is about twenty times the food/drink share in that trace
(0.0099). Nothing here is a claim about multi-block warehouses,
batching, or picker routing. Section 2.1 stated that restriction before
any number. The method name DAHS is historical; Section 3 does not model
disruptions.

\subsection{No online adaptation after
deployment}\label{no-online-adaptation-after-deployment}

The ranker is frozen. Regime posteriors are computed from the fitted
mixture, not updated from live outcomes. A shift that drifts away from
the training mixture gets the wrong specialist. The misspecification
grid (Section 6.11) is the evidence we have that this matters, and
Always-COVERT's shallower slope is a reminder that a dumb static rule
does not carry a wrong model.

Further limitations, stated so they are not rediscovered as objections.

\emph{Partial observability.} The aliasing rate among \(\phi\)-identical
training states that disagree on the cost-minimising rule is \(0.145\)
on 1,927 pairs, with mean regret \(0.30\) and an estimated \(3.7\%\)
share of the achievable benefit available from rule selection. That
number is an upper bound on the part of the residual regret attributable
to \(\phi\), not a reason to treat \(\phi\) as a state.

\emph{Regime order.} \(K^\star=12\) is the edge of the swept grid. BIC
is still falling (\(-168{,}901\) at \(K=8\), \(-222{,}525\) at \(K=10\),
\(-240{,}139\) at \(K=12\)). Mean ARI is \(0.970\), so the clustering is
reproducible; the model order is not identified. We did not widen the
grid, because doing so would retrain Stage 3 onward.

\emph{A pool of six is what a dispatcher already has.} Nine candidates
were considered and six retained (Section 3.6). Adding generated rules
would change the labelling budget and the action-coverage comparison,
and we did not do it.

\emph{The FQI comparison is a 4\% cost gap after a logger correction.}
It is not a demonstration that bootstrapped value learning cannot close
that gap with a different model class or a different behaviour policy.

\begin{center}\rule{0.5\linewidth}{0.5pt}\end{center}

\section{Conclusion}\label{conclusion}

We studied a selection hyper-heuristic for periodic-review warehouse
dispatching under two deadline clocks, trained by offline truncated
rollouts of a screened rule pool, on one warehouse simulator. The
training mechanism is not new. The question was whether the distilled
ranker recovers the online teacher on that simulator.

It does not. On the default corpus, online truncated lookahead is
cheaper (\(357\)--\(363\)) than the distilled ranker (\(382\)).
Always-COVERT is the static champion (\(455\); \(49/50\) shifts). Fitted
Q is about 4\% worse; that gap is not a result about value learning. PPO
is not a matched simulation-step budget; untuned PPO costs \(612\), or
\(450\) after observation and reward normalisation. Single-sample labels
match the deployed \(M=20\) configuration on composite cost. A one-step
(\(\tau=1\)) label is a null on mean \(J\) under the paired interval. A
two-step label is cheaper than deployed \(\tau=4\) on the mean
(\(-9.11\), \([-17.61,-1.35]\); 21/20/9, median 0), with no median
support. On the 12-cell robustness grid DAHS wins none of twelve cells
among five frozen methods. Distillation is a \(4.24\) ms decision that
spends \(5\)--\(7\%\) more than a \(670\) ms teacher; that \(670\) ms is
\(0.07\%\) of a 15-minute epoch.

The most useful extension is the one the dynamic-dispatching literature
already runs: replace hard truncation with a learned value tail (Ulmer
et al. 2019; Goodson et al. 2017), and let the horizon be
state-dependent. That construction would likely permit a shorter
\(\tau\), which is attractive here because Table 8 says \(\tau=2\)
already beats deployed \(\tau=4\) on mean \(J\). Two other extensions
follow from the campaign: widen the regime grid until \(K^\star\) is
interior, or drop the mixture; and test the null on a second
environment. One simulator cannot support a claim about training signals
in general.

\begin{center}\rule{0.5\linewidth}{0.5pt}\end{center}

\section{Appendix A. State features, their provenance, and their
redundancy}\label{appendix-a.-state-features-their-provenance-and-their-redundancy}

\subsection{Where each feature comes
from}\label{where-each-feature-comes-from}

This appendix lists every feature with its group and the source or
design rationale. The table is generated directly from
\texttt{simulation.state\_extractor.FEATURE\_PROVENANCE} by
\texttt{experiments/feature\_analysis.py}, so the manuscript and the
deployed feature map cannot drift apart.

The observation is \(\phi(S_t) \in \mathbb{R}^{26}\). It is an
observation, not a state (Section 3.2).

{\def\LTcaptype{none} % do not increment counter
\begin{longtable}[]{@{}
  >{\raggedright\arraybackslash}p{(\linewidth - 2\tabcolsep) * \real{0.5000}}
  >{\raggedright\arraybackslash}p{(\linewidth - 2\tabcolsep) * \real{0.5000}}@{}}
\toprule\noalign{}
\begin{minipage}[b]{\linewidth}\raggedright
Feature
\end{minipage} & \begin{minipage}[b]{\linewidth}\raggedright
Rationale
\end{minipage} \\
\midrule\noalign{}
\endhead
\bottomrule\noalign{}
\endlastfoot
\texttt{queue\_length} & Standard congestion state; de Koster et
al.~(2007). \\
\texttt{mean\_queue\_age} & Waiting-time proxy; distinguishes fresh from
stale backlog. \\
\texttt{max\_queue\_age} & Tail of the waiting-time distribution ---
starvation detector. \\
\texttt{pct\_critical} & Share of queue within 30 min of its due
time. \\
\texttt{pct\_perishable} & Share of the queue that carries a product
clock. The expiry-rule mask reads this; on the screened pool the mask is
a no-op. \\
\texttt{n\_arrivals\_last\_interval} & Short-run demand shock; wave
arrivals, Boysen et al.~(2019). \\
\texttt{labor\_utilization} & Classical queueing load indicator rho. \\
\texttt{n\_pickers\_busy} & Absolute capacity remaining this epoch. \\
\texttt{mean\_pickup\_time\_recent} & Realised service-rate estimate;
drifts with order mix. \\
\texttt{n\_orders\_late\_so\_far} & Realised failures to date; regime
indicator. \\
\texttt{n\_orders\_at\_risk\_30min} & Count with negative slack inside
30 min --- the actionable set. \\
\texttt{mean\_slack\_minutes} & Mean d - t - p; the ATC/MS/COVERT
decision variable. \\
\texttt{std\_slack\_minutes} & Slack dispersion; separates uniform from
bimodal pressure. \\
\texttt{mean\_processing\_time\_remaining} & Expected work content of
the queue. \\
\texttt{pct\_high\_priority} & Share of the economically weighted
tail. \\
\texttt{pct\_expiring\_30min} & Product-clock analogue of
pct\_critical. \\
\texttt{mean\_expiry\_slack} & Mean x - t - p over perishables; FEFO's
decision variable. \\
\texttt{n\_spoiled\_so\_far} & Finished spoilage to date: orders with
\texttt{finish\_time} \(\le t\) that are spoiled. Queued expired
perishables are invisible until picked (Section 3.3 keeps them in the
queue). \\
\texttt{arrival\_rate\_recent\_60min} & Non-stationarity detector;
empirical lambda-hat. \\
\texttt{queue\_length\_lag\_1} & Congestion trend; standard lag
structure. \\
\texttt{queue\_length\_lag\_2} & Congestion trend. \\
\texttt{queue\_length\_lag\_3} & Congestion trend. \\
\texttt{failure\_rate\_lag\_1} & Failure trend. \\
\texttt{failure\_rate\_lag\_2} & Failure trend. \\
\texttt{failure\_rate\_lag\_3} & Failure trend. \\
\texttt{interval\_index\_in\_shift} & Position in the finite horizon;
end-of-shift effects. \\
\end{longtable}
}

\subsection{Degenerate features that were
dropped}\label{degenerate-features-that-were-dropped}

\textbf{Two features are omitted} from the deployed map.
\texttt{time\_to\_next\_expected\_carrier} was computed as \(1/\lambda\)
and was therefore \emph{constant} within any one configuration: zero
variance, zero information. \texttt{intervals\_remaining} was an exact
affine function of \texttt{interval\_index\_in\_shift}, the two summing
to \(N = 32\) by construction. Together they made the feature matrix
exactly singular. That is not a cosmetic defect: the regime layer
(Section 4.4) fits a full-covariance Gaussian mixture on these columns,
so the covariance was singular, only the ridge term \texttt{reg\_covar}
prevented a failure, and each additional mixture component bought
likelihood by collapsing further onto the degenerate directions.

\textbf{Three features record the product clock}:
\texttt{pct\_expiring\_30min}, \texttt{mean\_expiry\_slack} and
\texttt{n\_spoiled\_so\_far}. The product deadline enters the objective
(Section 3.3), and a selector cannot act on a constraint it cannot
observe. FEFO's decision variable, expiry slack, has to be visible to
the ranker that decides when to deploy FEFO.

\textbf{One feature is named for the primary metric.} The failure-rate
lags are \texttt{failure\_rate\_lag\_1..3}.

\subsection{Redundancy analysis}\label{redundancy-analysis}

\texttt{experiments/feature\_analysis.py} runs four diagnostics on the
training feature matrix: near-constant columns; Pearson and Spearman
correlation with pairs above \(|r| = 0.95\) flagged; variance inflation
factors, obtained by regressing each feature on the remaining ones; and
correlation-distance clustering with one nominated representative per
cluster.

No column is near-constant. Ten pairs exceed \(|r|=0.95\), of which
\texttt{labor\_utilization} with \texttt{n\_pickers\_busy} is \(r=1\)
(they encode the same capacity residual) and the failure-rate lags with
\texttt{n\_orders\_late\_so\_far} are \(r \ge 0.979\). Fourteen features
have VIF \(\ge 10\): \texttt{queue\_length},
\texttt{n\_arrivals\_last\_interval}, \texttt{n\_pickers\_busy},
\texttt{labor\_utilization}, the three queue-length lags, the three
failure-rate lags, \texttt{n\_orders\_late\_so\_far},
\texttt{mean\_queue\_age}, \texttt{max\_queue\_age}, and
\texttt{mean\_slack\_minutes}. The clustering nominates dropping
\texttt{labor\_utilization}, \texttt{mean\_queue\_age},
\texttt{queue\_length\_lag\_1}, \texttt{queue\_length\_lag\_2}, and the
three failure-rate lags. We did not drop them from the deployed model:
the \texttt{top5\_features} ablation (Section 6.8) already asks whether
the long tail of the map earns its keep, and it does not, at least not
detectably on 50 test shifts. The VIF column is a warning that
coefficient-level stories about individual features are not identified,
which is why Section 6.6 reports SHAP on the fitted trees rather than
linear weights.

After dropping the two degenerate features of Section A.2, the regime
BIC still selects \(K^\star=12\) at the edge of
\(\{2,3,4,5,6,7,8,10,12\}\) (mean ARI \(0.970\), above the \(0.85\)
threshold). BIC values fall from \(+36{,}871\) at \(K=2\) to
\(-240{,}139\) at \(K=12\) and do not turn. The grid chose \(K\), not a
mode in the data. Input dimension is \(38\).

\section{Appendix B. Configuration and
hyperparameters}\label{appendix-b.-configuration-and-hyperparameters}

Every value below is read from \texttt{config.yaml} in the accompanying
repository, which is the source the campaign shares, except where a grid
winner is named.

\textbf{Fitted Q-iteration.} \texttt{config.yaml} default trees are
\texttt{max\_depth} = 4, \texttt{n\_estimators} = 200,
\texttt{learning\_rate} = 0.05, \(\gamma=0.99\), 20 iterations. The E9
grid winner written to \texttt{results/E9/hp\_winner.json} is
\texttt{max\_depth} = 4, \texttt{n\_estimators} = 500,
\texttt{learning\_rate} = 0.05, \(\gamma=0.9\). Table 5 uses the winner.

\textbf{Simulator.} 8-hour shift; \(N = 32\) review intervals of
\(L = 15\) minutes; \(m = 10\) pickers; queue capacity 200 orders, with
overflow recorded as rejected demand rather than discarded. Arrivals are
Poisson at a nominal 1.65 orders/minute. Order attributes: processing
time Triangular\((2, 5, 12)\) min; customer window \(d_o - a_o\)
Triangular\((15, 45, 90)\) min; shelf life \(x_o - a_o\)
Triangular\((20, 60, 120)\) min for the 20\% of orders that are
perishable, with \(x_o = \infty\) otherwise; priority classes \(\{\)low,
medium, high\(\}\) drawn at \((0.50, 0.35, 0.15)\) with economic weights
\(w_o \in \{1, 2, 4\}\). Section 3.4 gives the provenance of each.

\textbf{Objective} (Section 3.3). \(W_b = 3.0\) per late shipment,
\(W_t = 0.2\) per minute of lateness, \(W_s = 5.0\) per spoiled
perishable, \(W_h = 0.005\) per order still queued at each potential
evaluation (including window labels at \(t+\tau L\)); shift-level \(J\)
uses \(|Q_T|\). Every per-order charge is multiplied by \(w_o\). An
unweighted objective with no \(W_s\) term would grade WSPT and ATC
against a criterion they were not designed for, and would price an
abandoned order at \(0.005\) against \(3.0\) for one served late.

\textbf{Shift corpora.} Seeds drawn from one \texttt{SeedSequence} (root
42) and partitioned into three disjoint contiguous blocks: 250 training,
30 calibration, 50 test. The calibration block exists so that ATC's and
COVERT's look-ahead scales can be fitted without touching the training
or test shifts.

\textbf{Labelling} (Section 4.3). Rollout horizon \(\tau = 4\);
\(M = 20\) independent continuations per state-rule cell under common
random numbers, with the per-cell standard error recorded alongside
every label; behaviour policy \texttt{random} --- not round robin,
because the interval index is itself an observed feature, which would
make round robin a deterministic function of the state (Section 6.10).
Tempered softmax with a \textbf{per-row} temperature
\(T(s) = \beta\,\sigma(s)\), where \(\sigma(s)\) is the standard
deviation of that state's own cost vector and \(\beta\) is searched so
that the median training-label entropy falls within
\([0.216, 0.505] \times \log|\mathcal{H}|\). Campaign fit:
\(\beta=0.470\), median entropy \(0.638\) nats, band \([0.387, 0.905]\).
Expiry-rule mask threshold 0.05 on the perishable fraction (a no-op on
the screened pool, which contains no expiry-only rule). Test-corpus
ambiguity filter at \(\theta = 2.2/|\mathcal{H}|\), never applied to the
training corpus. The hard-label ablation of Section 6.8 replaces the
tempered softmax with the one-hot arg-max of the same cost vector and is
otherwise identical.

\textbf{Rule pool} (Section 3.6). Nine candidates before screening:
FIFO, EDD, EEDD, FEFO, WSPT, ATC, MS, MDD, COVERT. Stage 1 retained
\texttt{{[}EEDD,\ COVERT,\ MS,\ ATC,\ MDD,\ EDD{]}} and fitted ATC
\(k^\star=3.0\) (portfolio; standalone \(1.5\)) and COVERT
\(k^\star=4.0\).

\textbf{Regime layer} (Section 4.4). Gaussian mixture, full covariance,
with \(K\) selected by BIC over \(K \in \{2,3,4,5,6,7,8,10,12\}\) and 5
EM restarts per \(K\); stability checked by the mean adjusted Rand index
over 10 refits against a 0.85 threshold. \(K^\star=12\) at the grid
edge, mean ARI \(0.970\).

\textbf{Ranker.} Gradient-boosted trees (Chen and Guestrin 2016),
\texttt{multi:softprob} objective, sample-weighted by inverse label
entropy. Hyperparameters selected from an 18-configuration grid
(\texttt{max\_depth} \(\in \{4,6,8\}\) × \texttt{n\_estimators} \(\in
\{200,500,1000\}\) × \texttt{learning\_rate} \(\in \{0.03,0.1\}\)) by
5-fold cross-validation grouped on \texttt{shift\_id}; isotonic
calibration on a 20\% held-out shift split. Selected:
\texttt{max\_depth} = 6, \texttt{n\_estimators} = 200,
\texttt{learning\_rate} = 0.03. Test ECE after isotonic: \(0.0213\) (pre
\(0.1700\)).

\textbf{Switching controller} (Section 4.6). Minimum dwell
\(T_{\min} = 2\) intervals; entropy gate at half the maximum entropy.

\section{Appendix C. Fitting the input distributions to a real order
trace}\label{appendix-c.-fitting-the-input-distributions-to-a-real-order-trace}

\textbf{Source.} Olist Brazilian e-commerce public dataset (Olist 2018),
approximately 100k orders. Inter-arrival times are computed as
within-day differences of order timestamps, which removes the multi-day
growth trend that otherwise dominates the series. Because the trace is
measured in days and the simulator in minutes, all comparisons are on
mean-normalised samples: the trace fixes distribution \emph{shape}, and
the operating rate is set separately.

Candidate families for inter-arrival times: exponential, lognormal,
gamma, Weibull. Selected: \textbf{lognormal} (\(n=98{,}241\), AIC
\(161{,}984\)). Exponential is worse by \(\Delta\)AIC \(=34{,}502\).
Shape tests against the simulator's Poisson default: KS \(D=0.153\),
Wasserstein-1 \(0.443\); simulated CV \(1.00\) versus real \(2.68\).
Poisson is therefore an operating point, not a fitted arrival process.
Section 6.7's Olist-bootstrap replay is the robustness check that uses
the empirical burstiness.

Customer windows: the AIC-best family on the Olist
purchase-to-estimated-delivery sample is \textbf{lognormal}. The
campaign does not deploy that MLE. It uses triangular ((15, 45, 90)) as
a warehouse-scale operating-point envelope. Closest shape match of the
three inputs against that envelope (KS \(D=0.039\), subsampled
\(p=0.022\), \(W_1=0.035\)).

Processing time: not fitted. The trace field is purchase-to-approval,
not pick time. The triangular ((2, 5, 12)) is a literature three-point
standard (Tompkins et al. 2010). Shape test against that proxy is
correspondingly poor (KS \(D=0.686\)).

Perishable fraction: design parameter \(0.20\) against Olist food/drink
share \(0.0099\). Shelf life has no public analogue.

\begin{center}\rule{0.5\linewidth}{0.5pt}\end{center}

\section*{Data availability}

The simulator, training pipeline, configuration and result artifacts
that produced the live tables in this paper are in the accompanying
repository. Table 1 is the exception: it used the eight-rule
\texttt{SCREENING\_POOL} (FIFO omitted) on 30 calibration shifts, stored
at \texttt{results/S1\_perishability/}. The live \texttt{resolve\_pool}
in \texttt{config.yaml} is the six-rule deployed pool, so reproducing
Table 1 from the current config will not recover it. CAOR requires that
the data and code underlying the work be deposited (Guide for Authors,
Option C). A persistent archive with a DOI will be linked at acceptance;
until then the versioned source is
\url{https://github.com/Vittal-Mukunda/Disruption-Aware-Heuristic-Scheduling}.

\section*{CRediT authorship contribution statement}

\textbf{Vittal Mukunda:} Conceptualization, Methodology, Software,
Formal analysis, Writing --- original draft. \textbf{Atharva Somani:}
Software, Investigation, Validation, Data curation. \textbf{Pranjal
Malaiya:} Methodology, Writing --- review and editing, Supervision.

\section*{Declaration of competing interest}

The authors declare that they have no known competing financial
interests or personal relationships that could have appeared to
influence the work reported in this paper.

\section*{Funding}

This research did not receive any specific grant from funding agencies
in the public, commercial, or not-for-profit sectors.

\section*{Acknowledgements}

We thank the editors and the anonymous referees for comments on an
earlier version of this work.

\protect\phantomsection\label{refs}
\begin{CSLReferences}{1}{1}
\bibitem[\citeproctext]{ref-baker1982mdd}
Baker, Kenneth R., and J. Will M. Bertrand. 1982. {``A Dynamic Priority
Rule for Scheduling Against Due-Dates.''} \emph{Journal of Operations
Management} 3 (1): 37--42.
\url{https://doi.org/10.1016/0272-6963(82)90020-1}.

\bibitem[\citeproctext]{ref-bertsekas2020rollout}
Bertsekas, Dimitri P. 2020. \emph{Rollout, Policy Iteration, and
Distributed Reinforcement Learning}. Athena Scientific.

\bibitem[\citeproctext]{ref-boysen2025warehousing}
Boysen, Nils, and René de Koster. 2025. {``50 Years of Warehousing
Research---an Operations Research Perspective.''} \emph{European Journal
of Operational Research} 320 (3): 449--64.
\url{https://doi.org/10.1016/j.ejor.2024.03.026}.

\bibitem[\citeproctext]{ref-boysen2019warehousing}
Boysen, Nils, René de Koster, and Felix Weidinger. 2019. {``Warehousing
in the e-Commerce Era: A Survey.''} \emph{European Journal of
Operational Research} 277 (2): 396--411.
\url{https://doi.org/10.1016/j.ejor.2018.08.023}.

\bibitem[\citeproctext]{ref-branke2016automated}
Branke, Jürgen, Su Nguyen, Christoph W. Pickardt, and Mengjie Zhang.
2016. {``Automated Design of Production Scheduling Heuristics: A
Review.''} \emph{IEEE Transactions on Evolutionary Computation} 20 (1):
110--24. \url{https://doi.org/10.1109/TEVC.2015.2429314}.

\bibitem[\citeproctext]{ref-carroll1965covert}
Carroll, Donald C. 1965. {``Heuristic Sequencing of Single and Multiple
Component Jobs.''} PhD thesis, Massachusetts Institute of Technology.

\bibitem[\citeproctext]{ref-chen2016xgboost}
Chen, Tianqi, and Carlos Guestrin. 2016. {``{XGBoost}: A Scalable Tree
Boosting System.''} \emph{Proceedings of the 22nd ACM SIGKDD
International Conference on Knowledge Discovery and Data Mining},
785--94. \url{https://doi.org/10.1145/2939672.2939785}.

\bibitem[\citeproctext]{ref-zhang2024lstmppo}
Chen, Wei, Zequn Zhang, Dunbing Tang, et al. 2024. {``Probing an
{LSTM}-{PPO}-Based Reinforcement Learning Algorithm to Solve Dynamic Job
Shop Scheduling Problem.''} \emph{Computers \& Industrial Engineering}
197: 110633. \url{https://doi.org/10.1016/j.cie.2024.110633}.

\bibitem[\citeproctext]{ref-cheng2024drlhyperheuristic}
Cheng, Bayi, Lingjun Wang, Qi Tan, and Mi Zhou. 2024. {``A Deep
Reinforcement Learning Hyper-Heuristic to Solve Order Batching Problem
with Mobile Robots.''} \emph{Applied Intelligence} 54: 6865--87.
\url{https://doi.org/10.1007/s10489-024-05532-9}.

\bibitem[\citeproctext]{ref-conway1967theory}
Conway, Richard W., William L. Maxwell, and Louis W. Miller. 1967.
\emph{Theory of Scheduling}. Addison-Wesley.

\bibitem[\citeproctext]{ref-dimitrakakis2008rollout}
Dimitrakakis, Christos, and Michail G. Lagoudakis. 2008. {``Rollout
Sampling Approximate Policy Iteration.''} \emph{Machine Learning} 72
(3): 157--71. \url{https://doi.org/10.1007/s10994-008-5069-3}.

\bibitem[\citeproctext]{ref-dokeroglu2024hyperheuristics}
Dokeroglu, Tansel, Tayfun Kucukyilmaz, and El-Ghazali Talbi. 2024.
{``Hyper-Heuristics: A Survey and Taxonomy.''} \emph{Computers \&
Industrial Engineering} 187: 109815.
\url{https://doi.org/10.1016/j.cie.2023.109815}.

\bibitem[\citeproctext]{ref-drake2020hyperheuristics}
Drake, John H., Ahmed Kheiri, Ender Özcan, and Edmund K. Burke. 2020.
{``Recent Advances in Selection Hyper-Heuristics.''} \emph{European
Journal of Operational Research} 285 (2): 405--28.
\url{https://doi.org/10.1016/j.ejor.2019.07.073}.

\bibitem[\citeproctext]{ref-durasevic2022dispatching}
Durašević, Marko, and Domagoj Jakobović. 2022. {``Selection of
Dispatching Rules Evolved by Genetic Programming in Dynamic Unrelated
Machines Scheduling Based on Problem Characteristics.''} \emph{Journal
of Computational Science} 61: 101649.
\url{https://doi.org/10.1016/j.jocs.2022.101649}.

\bibitem[\citeproctext]{ref-ernst2005fqi}
Ernst, Damien, Pierre Geurts, and Louis Wehenkel. 2005. {``Tree-Based
Batch Mode Reinforcement Learning.''} \emph{Journal of Machine Learning
Research} 6: 503--56.

\bibitem[\citeproctext]{ref-farahmand2015capi}
Farahmand, Amir-massoud, Doina Precup, André M. S. Barreto, and Mohammad
Ghavamzadeh. 2015. {``Classification-Based Approximate Policy
Iteration.''} \emph{IEEE Transactions on Automatic Control} 60 (11):
2989--93. \url{https://doi.org/10.1109/TAC.2015.2418411}.

\bibitem[\citeproctext]{ref-fern2006api}
Fern, Alan, Sungwook Yoon, and Robert Givan. 2006. {``Approximate Policy
Iteration with a Policy Language Bias: Solving Relational {M}arkov
Decision Processes.''} \emph{Journal of Artificial Intelligence
Research} 25: 75--118. \url{https://doi.org/10.1613/jair.1700}.

\bibitem[\citeproctext]{ref-geng2016ldl}
Geng, Xin. 2016. {``Label Distribution Learning.''} \emph{IEEE
Transactions on Knowledge and Data Engineering} 28 (7): 1734--48.
\url{https://doi.org/10.1109/TKDE.2016.2545658}.

\bibitem[\citeproctext]{ref-vangils2018designing}
Gils, Teun van, Katrien Ramaekers, An Caris, and René B. M. de Koster.
2018. {``Designing Efficient Order Picking Systems by Combining Planning
Problems: State-of-the-Art Classification and Review.''} \emph{European
Journal of Operational Research} 267 (1): 1--15.
\url{https://doi.org/10.1016/j.ejor.2017.09.002}.

\bibitem[\citeproctext]{ref-goodson2016restocking}
Goodson, Justin C., Barrett W. Thomas, and Jeffrey W. Ohlmann. 2016.
{``Restocking-Based Rollout Policies for the Vehicle Routing Problem
with Stochastic Demand and Duration Limits.''} \emph{Transportation
Science} 50 (2): 591--607. \url{https://doi.org/10.1287/trsc.2015.0591}.

\bibitem[\citeproctext]{ref-goodson2017rolloutframework}
Goodson, Justin C., Barrett W. Thomas, and Jeffrey W. Ohlmann. 2017.
{``A Rollout Algorithm Framework for Heuristic Solutions to
Finite-Horizon Stochastic Dynamic Programs.''} \emph{European Journal of
Operational Research} 258 (1): 216--29.
\url{https://doi.org/10.1016/j.ejor.2016.09.040}.

\bibitem[\citeproctext]{ref-hanjung2025imitation}
Han, Young-Suk, and Jae-Yoon Jung. 2025. {``Imitation Learning of Job
Dispatching Data Combining Classification and Ranking.''}
\emph{Computers \& Industrial Engineering} 207: 111305.
\url{https://doi.org/10.1016/j.cie.2025.111305}.

\bibitem[\citeproctext]{ref-jackson1955edd}
Jackson, James R. 1955. \emph{Scheduling a Production Line to Minimize
Maximum Tardiness}. Research Report 43. Management Sciences Research
Project, University of California, Los Angeles.

\bibitem[\citeproctext]{ref-klapp2018dispatchwaves}
Klapp, Mathias A., Alan L. Erera, and Alejandro Toriello. 2018a. {``The
Dynamic Dispatch Waves Problem for Same-Day Delivery.''} \emph{European
Journal of Operational Research} 271 (2): 519--34.
\url{https://doi.org/10.1016/j.ejor.2018.05.032}.

\bibitem[\citeproctext]{ref-klapp2018onedim}
Klapp, Mathias A., Alan L. Erera, and Alejandro Toriello. 2018b. {``The
One-Dimensional Dynamic Dispatch Waves Problem.''} \emph{Transportation
Science} 52 (2): 402--15. \url{https://doi.org/10.1287/trsc.2016.0682}.

\bibitem[\citeproctext]{ref-dekoster2007orderpicking}
Koster, René de, Tho Le-Duc, and Kees Jan Roodbergen. 2007. {``Design
and Control of Warehouse Order Picking: A Literature Review.''}
\emph{European Journal of Operational Research} 182 (2): 481--501.

\bibitem[\citeproctext]{ref-lagoudakis2003rcpi}
Lagoudakis, Michail G., and Ronald Parr. 2003. {``Reinforcement Learning
as Classification: Leveraging Modern Classifiers.''} \emph{Proceedings
of the 20th International Conference on Machine Learning (ICML)},
424--31.

\bibitem[\citeproctext]{ref-law2000simulation}
Law, Averill M., and W. David Kelton. 2000. \emph{Simulation Modeling
and Analysis}. 3rd ed. McGraw-Hill.

\bibitem[\citeproctext]{ref-li2010linucb}
Li, Lihong, Wei Chu, John Langford, and Robert E. Schapire. 2010. {``A
Contextual-Bandit Approach to Personalized News Article
Recommendation.''} \emph{Proceedings of the 19th International
Conference on World Wide Web (WWW)}, 661--70.
\url{https://doi.org/10.1145/1772690.1772758}.

\bibitem[\citeproctext]{ref-liu2025mlscheduling}
Liu, Anbang, Shaochong Lin, Jingchuan Chen, Peng Wu, and Zuojun Max
Shen. 2025. \emph{Machine Learning for Scheduling: A Paradigm Shift from
Solver-Centric to Data-Centric Approaches}.
\url{https://arxiv.org/abs/2512.22642}.

\bibitem[\citeproctext]{ref-lundberg2017shap}
Lundberg, Scott M., and Su-In Lee. 2017. {``A Unified Approach to
Interpreting Model Predictions.''} \emph{Advances in Neural Information
Processing Systems} 30.

\bibitem[\citeproctext]{ref-mahmoudinazlou2025drl}
Mahmoudinazlou, Sasan, Abhay Sobhanan, Hadi Charkhgard, Ali Eshragh, and
George Dunn. 2025. {``Deep Reinforcement Learning for Dynamic Order
Picking in Warehouse Operations.''} \emph{Computers \& Operations
Research} 182: 107112. \url{https://doi.org/10.1016/j.cor.2025.107112}.

\bibitem[\citeproctext]{ref-mouelhi2010neural}
Mouelhi-Chibani, Wiem, and Henri Pierreval. 2010. {``Training a Neural
Network to Select Dispatching Rules in Real Time.''} \emph{Computers \&
Industrial Engineering} 58 (2): 249--56.
\url{https://doi.org/10.1016/j.cie.2009.03.008}.

\bibitem[\citeproctext]{ref-nguyen2017gpsurvey}
Nguyen, Su, Yi Mei, and Mengjie Zhang. 2017. {``Genetic Programming for
Production Scheduling: A Survey with a Unified Framework.''}
\emph{Complex \& Intelligent Systems} 3 (1): 41--66.
\url{https://doi.org/10.1007/s40747-017-0036-x}.

\bibitem[\citeproctext]{ref-olist2018dataset}
Olist. 2018. \emph{Brazilian e-Commerce Public Dataset by Olist}.
Kaggle.

\bibitem[\citeproctext]{ref-pinedo2016scheduling}
Pinedo, Michael L. 2016. \emph{Scheduling: Theory, Algorithms, and
Systems}. 5th ed. Springer.
\url{https://doi.org/10.1007/978-3-319-26580-3}.

\bibitem[\citeproctext]{ref-powell2019unified}
Powell, Warren B. 2019. {``A Unified Framework for Stochastic
Optimization.''} \emph{European Journal of Operational Research} 275
(3): 795--821. \url{https://doi.org/10.1016/j.ejor.2018.07.014}.

\bibitem[\citeproctext]{ref-powell2022rlso}
Powell, Warren B. 2022. \emph{Reinforcement Learning and Stochastic
Optimization: A Unified Framework for Sequential Decisions}. Wiley.

\bibitem[\citeproctext]{ref-vanremmerden2025offlineld}
Remmerden, Jesse van, Zaharah Bukhsh, and Yingqian Zhang. 2025.
{``Offline Reinforcement Learning for Learning to Dispatch for Job Shop
Scheduling.''} \emph{Machine Learning} 114: 191.
\url{https://doi.org/10.1007/s10994-025-06826-w}.

\bibitem[\citeproctext]{ref-roodbergen2006layout}
Roodbergen, Kees Jan, and Iris F. A. Vis. 2006. {``A Model for Warehouse
Layout.''} \emph{IIE Transactions} 38 (10): 799--811.

\bibitem[\citeproctext]{ref-sargent2013vandv}
Sargent, Robert G. 2013. {``Verification and Validation of Simulation
Models.''} \emph{Proceedings of the Winter Simulation Conference}.

\bibitem[\citeproctext]{ref-schulman2017ppo}
Schulman, John, Filip Wolski, Prafulla Dhariwal, Alec Radford, and Oleg
Klimov. 2017. \emph{Proximal Policy Optimization Algorithms}.
\url{https://arxiv.org/abs/1707.06347}.

\bibitem[\citeproctext]{ref-shiue2020rl}
Shiue, Yeou-Ren, Ken-Chuan Lee, and Chao-Ton Su. 2020. {``A
Reinforcement Learning Approach to Dynamic Scheduling in a Product-Mix
Flexibility Environment.''} \emph{IEEE Access} 8: 106542--53.
\url{https://doi.org/10.1109/ACCESS.2020.3000781}.

\bibitem[\citeproctext]{ref-smith1956wspt}
Smith, Wayne E. 1956. {``Various Optimizers for Single-Stage
Production.''} \emph{Naval Research Logistics Quarterly} 3 (1-2):
59--66. \url{https://doi.org/10.1002/nav.3800030106}.

\bibitem[\citeproctext]{ref-stockermann2025drlscalability}
Stöckermann, Patrick, Henning Südfeld, Alessandro Immordino, et al.
2025. {``Scalability of Reinforcement Learning Methods for Dispatching
in Semiconductor Frontend Fabs: A Comparison of Open-Source Models with
Real Industry Datasets.''} \emph{The International Journal of Advanced
Manufacturing Technology}, ahead of print.
\url{https://doi.org/10.1007/s00170-025-16117-2}.

\bibitem[\citeproctext]{ref-tassel2023rljssp}
Tassel, Pierre, Martin Gebser, and Konstantin Schekotihin. 2023. {``An
End-to-End Reinforcement Learning Approach for Job-Shop Scheduling
Problems Based on Constraint Programming.''} \emph{Proceedings of the
International Conference on Automated Planning and Scheduling} 33 (1):
614--22. \url{https://doi.org/10.1609/icaps.v33i1.27243}.

\bibitem[\citeproctext]{ref-tompkins2010facilities}
Tompkins, James A., John A. White, Yavuz A. Bozer, and J. M. A.
Tanchoco. 2010. \emph{Facilities Planning}. 4th ed. John Wiley \& Sons.

\bibitem[\citeproctext]{ref-ulmer2019anticipation}
Ulmer, Marlin W. 2019. {``Anticipation Versus Reactive Reoptimization
for Dynamic Vehicle Routing with Stochastic Requests.''} \emph{Networks}
73 (3): 277--91. \url{https://doi.org/10.1002/net.21861}.

\bibitem[\citeproctext]{ref-ulmer2019offlineonline}
Ulmer, Marlin W., Justin C. Goodson, Dirk C. Mattfeld, and Marco Hennig.
2019. {``Offline--Online Approximate Dynamic Programming for Dynamic
Vehicle Routing with Stochastic Requests.''} \emph{Transportation
Science} 53 (1): 185--202. \url{https://doi.org/10.1287/trsc.2017.0767}.

\bibitem[\citeproctext]{ref-ulmer2020modeling}
Ulmer, Marlin W., Justin C. Goodson, Dirk C. Mattfeld, and Barrett W.
Thomas. 2020. {``On Modeling Stochastic Dynamic Vehicle Routing
Problems.''} \emph{EURO Journal on Transportation and Logistics} 9 (2):
100008. \url{https://doi.org/10.1016/j.ejtl.2020.100008}.

\bibitem[\citeproctext]{ref-ulmer2018budgeting}
Ulmer, Marlin W., Dirk C. Mattfeld, and Felix Köster. 2018. {``Budgeting
Time for Dynamic Vehicle Routing with Stochastic Customer Requests.''}
\emph{Transportation Science} 52 (1): 20--37.
\url{https://doi.org/10.1287/trsc.2016.0719}.

\bibitem[\citeproctext]{ref-vepsalainen1987atc}
Vepsäläinen, Ari P. J., and Thomas E. Morton. 1987. {``Priority Rules
for Job Shops with Weighted Tardiness Costs.''} \emph{Management
Science} 33 (8): 1035--47. \url{https://doi.org/10.1287/mnsc.33.8.1035}.

\bibitem[\citeproctext]{ref-winkelhaus2020logistics4}
Winkelhaus, Sven, and Eric H. Grosse. 2020. {``Logistics 4.0: A
Systematic Review Towards a New Logistics System.''} \emph{International
Journal of Production Research} 58 (1): 18--43.
\url{https://doi.org/10.1080/00207543.2019.1612964}.

\bibitem[\citeproctext]{ref-wu1988multipass}
Wu, S. David, and Richard A. Wysk. 1988. {``Multi-Pass Expert Control
System --- a Control/Scheduling Structure for Flexible Manufacturing
Cells.''} \emph{Journal of Manufacturing Systems} 7 (2): 107--20.
\url{https://doi.org/10.1016/0278-6125(88)90018-0}.

\end{CSLReferences}

\end{document}
